% =====================================================================
%  R. Nielsen, AANDSVIRKELIGHED (1870) and VIDERE OM AANDSVIRKELIGHED, FORANLEDIGET VED PASTOR V. BIRKEDALS OG PASTOR C. MØLLERS INDLÆG (1871)
%  In: For Ide og Virkelighed. Et Tidsskrift, udgivet af R. Nielsen, Bj. Bjørnson og Rud. Schmidt.  Aargang 1870. Andet Bind (Bind 3) and Aargang 1871. Andet Bind (Bind 5).
%  Kjøbenhavn: Chr. Steen & Søns Forlag; Louis Kleins Bogtrykkeri -- Otto B. Wroblewskys Forlag; G. S. Wibes Bogtrykkeri.  Bind 3, pp. 289-320; Bind 5, pp. 1-38.  Antiqua.
%
%  TRANSCRIPTION -- GENERATED; DO NOT EDIT BY HAND.  Rasmus Nielsen's eight essays in the volumes 1870-1873 are built by ONE
%  pipeline (ebook-method/for-idee-1870-73/) into a combined working file, which split7.py cuts into seven transcriptions
%  (Aandsvirkelighed and its sequel, Videre om Aandsvirkelighed, are one: Hans, 2026-10-13).
%  Rebuild from the repo root:   sh ebook-method/for-idee-1870-73/rebuild.sh
%
%  STATUS: COMPLETE (2026-10-13): the volume's title page and the whole essay, Bind 3, pp. 289-320; Bind 5, pp. 1-38, with its notes.
%
%  EDITION.  The only printing.  Text of record: Nasjonalbiblioteket's scan (nb.no) of the volume, ~/bibliotek/Nielsen, Rasmus/1870-2-for-idee-og-virkelighed.pdf (NB item 5d1143e0fa49eff177d80b7a3a41fddd); ~/bibliotek/Nielsen, Rasmus/1871-2-for-idee-og-virkelighed.pdf (NB item 9bae6062efbd4219d60f56932f2a37a9) (SCANS.tsv),
%  page images fetched by nbfetch.py, 300 dpi colour, NO text layer; printed page N = PDF page N + 4.  No e-book, no
%  Wikisource text, no reprint, no earlier transcription found (checked 2026-10-13: the repo, the web, the Internet Archive,
%  da.wikisource; HathiTrust and the Google Books API could not be queried that day; Google Books holds at least the 1873
%  volume, nielsen/darwinismen's source -- a second copy, not used).
%
%  ACCURACY (2026-10-13; one collation for all eight texts, 231 pp., ~58,000 words).  The DRAFT is NB's own OCR (ALTO, ABBYY
%  FineReader 12), laid over the page images (combine.py) and turned into LaTeX by draft.py; it was collated word by word
%  with an independent tesseract reading of the same images (collate/check.py, witness 'ocr'): 1556 disagreements.
%  check.py's own rules settled 78.  Candidate rules were first scored BY SCRIPT on a text already read at the image
%  (calibrate.py: Et Synspunkt for Darwinismen, same volume as Verdensanskuelse), then measured on a BLIND PILOT of 152 crops
%  (reader A, 52/52 planted controls right); 17 rules passed (NB's glyph faults: o for ø, D/O/U for Ø, i for l, li for h,
%  l for I, opening » read «, words run together or broken apart; tesseract's: aa read an/ag/å, added accents; the
%  printer's space before a mark, house rule) and settled 772; the rest were read at the image BLIND (562 crops on 94
%  sheets, four readers, 188/188 planted controls right), ~160 decided in the main session from the readers' lines and ~40
%  at 300-dpi zooms.  STRUCTURE read at the image (reader K): every footnote whole (30; collated where they stand, then moved
%  to their marks, then compared word by word with the reading: notecmp.py, 9 differences settled at the image), the
%  headings, the section numerals, the essays' last lines and signatures, the three Greek words, the volume title pages.
%  LETTERSPACING (the print's emphasis, \emph): candidate lines by letter gaps and box width, calibrated on the words the
%  word readers had marked spaced (spaced.py): 361 lines read blind (61/62 roman controls), 25 doubtful extents read at a
%  zoom.  SWEEPS after OPEN 0: words run together that both OCRs share (lower-case + capital), the unknown-word sweep (678
%  words, ~20 suspects, 14 read at the image: 4 OCR faults, 7 misprints kept), --moved, the dropped-word list, the
%  W-shortens list of the blind decisions.
%  SAMPLED PAGE CHECK (pagecheck.py, seed 202610137): one page of each essay or part, pp. 2289, 2400, 3305, 4308, 5033,
%  5208, 6514, 8011 of the working numbering (Bind x 1000 + page), read whole and blind, 1965 words: 0 word errors,
%  1 punctuation (a compound's hyphen lost at a page turn: „Ny-|Holland“) and 1 letterspacing („to be or not to be“)
%  error, both mended.  95 % upper bound for word errors: 1.5 per 1000 words.
%  What this cannot see: a misprint both OCRs normalise into a real word.
%  The ø rule (Hans 2026-10-08): o-forms that are no word are set with ø; „stode“ (= stood) is a word and stays.
%  PRINTED AS IS Bind 5, p. 29: „Jevnf, Religionsfilosofien“ comma for the stop of Jevnf. (K M50, zoom z3)
%  PRINTED AS IS Bind 3, p. 317: „kvorpaa“ for hvorpaa (zoom z4)
%  Not transcribed: running heads, folios, sheet signatures („For Ide og Virkelighed. 1870. I. 20“), library stamps; the
%  large initial of each essay is set as a plain letter.  Quotation marks as printed (» « in the text, sometimes a closing
%  »; „ “ in some notes).  The notes are numbered ¹), ²) per page, as printed.  Greek as printed (textalpha).
%
%  TOOLS (ebook-method/for-idee-1870-73/): nbalto.py, combine.py, alto.py, ocr.py, draft.py, rules.py, calibrate.py,
%  sheets.py + reads.py + READER-BRIEF.md, spaced.py + spaced.tsv, structsheets.py, notecmp.py, unknown.py, zoom.py,
%  pagecheck.py + PAGECHECK-BRIEF.md, patches.py (PATCHES, PRENOTES, NOTES, POST, MISPRINTS), decisions.tsv, split7.py,
%  rebuild.sh.
% =====================================================================
\documentclass[12pt,a4paper,openany]{book}
\usepackage[utf8]{inputenc}
\usepackage[T1]{fontenc}
\usepackage{libertinus}
\usepackage{libertinust1math}
\usepackage{textalpha}
\usepackage[danish]{babel}
\usepackage[protrusion=true,expansion=true]{microtype}
\usepackage{setspace}
\usepackage[margin=1.3in]{geometry}
\usepackage{fancyhdr}
\usepackage{hyperref}
\hypersetup{pdftitle={Aandsvirkelighed}, pdfauthor={Rasmus Nielsen}, hidelinks}
\pagestyle{fancy}
\fancyhf{}
\fancyhead[LE,RO]{\thepage}
\fancyhead[C]{\textit{R. Nielsen: Aandsvirkelighed (1870--71)}}
\setstretch{1.3}
\usepackage{marginnote}
\renewcommand{\thefootnote}{\arabic{footnote})}   % the print numbers each page's notes ¹), ²)
\newcommand{\opage}[1]{\marginnote{\footnotesize\textit{[#1]}}}
% the essay's heading: title (and subtitle), centred, a short rule; the print's large initial is not reproduced
\newcommand{\essayhead}[2]{\begin{center}{\large #1}\ifx&#2&\else\\[0.3em]{\normalsize #2}\fi\\[0.5em]\rule{3em}{0.4pt}\end{center}\medskip\noindent}
\newcommand{\signature}[1]{\par\smallskip\hfill #1\hspace*{2em}\par}

\begin{document}

% --- front matter: the title page of Bind 3 (PDF 1 of 1870-2-for-idee-og-virkelighed.pdf), read at the image (reader K, m1);
% short wavy rules are set as plain rules; the library stamps are not transcribed.
\begin{titlepage}
\begin{center}
\vspace*{3em}
{\Huge For Ide og Virkelighed.}\\[1em]
\rule{8em}{0.4pt}\\[1em]
{\large Et Tidsskrift,}\\[0.8em]
{\small udgivet}\\[0.4em]
{\footnotesize af}\\[0.6em]
{\large R. Nielsen,}\\[0.4em]
{\large Bj. Bjørnson og Rud. Schmidt.}\\[1em]
\rule{8em}{0.4pt}\\[0.6em]
{\large Aargang 1870. Andet Bind.}\\[0.6em]
\rule{8em}{0.4pt}\\[4em]
Kjøbenhavn.\\[0.4em]
Chr. Steen \& Søns Forlag.\\[1.5em]
{\small Louis Kleins Bogtrykkeri.}
\end{center}
\end{titlepage}

% ===== PART I: Bind 3, pp. 289-320 =====
\essayhead{Aandsvirkelighed.}{}
% --- p. 289 ---
\setcounter{footnote}{0}%
\opage{289}Den aabenbarede Religion med Livets Ord, med Undere
og Tegn saavel i historiske som overhistoriske Kjendsgjerninger
er Aandens Værk. Her at ville gribe efter
det udvortes Virkelige og glemme Aanden, det evige
Indvortes, er ligesaa urimeligt som at ville gribe den
rene Aand udenfor de Kjendsgjerninger, i hvilke den
har aabenbaret sig. Livets og Sandhedens Aand maa
sees i Virkelighedens Lys, og Virkeligheden igjen i
Aandens Lys: Aabenbaringens Virkelighed er Aandsvirkelighed.

Som nu den Virkelighed er, der skal sees, saaledes
maa det Øie være, der skal see det Virkelige. For at
see de blot sandselige Ting er det nok at have et sandselig
sundt og skarpt Øie; i denne Henseende kan Dyret
endog have visse ensidige Fortrin for Mennesket. Hvad
der særlig udmærker det menneskelige Syn fremfor det
dyriske, er det Eiendommelige i Menneskets Blik, at
Mennesket seer tænkende paa Tingen.

Naturforskeren væbner sit Øie --- med Teleskopet,
for at see det Store, der ellers vilde være skjult i
uopnaaelig Fjernhed; med Mikroskopet, for at see, hvad
der trods den umiddelbare Nærhed ligger skjult i uen%
% --- p. 290 ---
\setcounter{footnote}{0}%
\opage{290}delig Lidenhed. Han vil see tilgavns, hvad der er at see
overalt i den synlige Verden; men han nøies ikke med
en sandselig Seen; Synet maa tjene hans Tanke; han
leder i alt det Synlige efter det Usynlige: han forsker
efter Love.

Tanken er aandig; en Verden af virkelige Love er
en aandig Verden. Vi kunne med fuld Ret sige, at
Naturforskeren i det Virkelige seer det Aandige, at
Virkeligheden staaer for hans Blik som en aandig Virkelighed;
men vi kunne ikke sige, at denne aandige
Virkelighed endnu er Aandsvirkelighed. Og hvorfor
ikke? Det er ikke Virkelighedens Skyld; nei, det er
Øiets Skyld eller, seet fra en anden Side, det er Belysningens
Skyld. Det aandige Lys, hvori Forskeren
alene kan, vil og bør see det hele virkelige Naturrige,
er Tankens, Fornuftens, Videnskabens Lys; men dette
Tankelys er i al sin Klarhed for almindeligt, koldt og
ensidigt til at være Aandens Lys.

At der i Virkelighedens frugtbare Egne endnu er
Andet at opdage end hvad der kan sees ved Videnskabens
almindelige Tankelys, derom afgive Konst og
Poesi da ogsaa Vidnesbyrd nok. Landskabsmaleren f.
Ex. kan aldrig blive træt af at betragte Naturen, thi
jo længere han betragter den, jo mere seer han i
den; den meddeler sig til ham, den betroer ham sine
Hemmeligheder: han gjør bestandig nye Opdagelser.
Naar Beskueren seer paa det malede Landskab og finder
Billedet vidunderlig skjønt, saa fristes han let til at
sige ved sig selv: dette er ingen tro Gjengivelse af
det Virkelige; det er Naturen, ikke som den virkelig
er, men som den i Ideens og Skjønhedens Navn bør
% --- p. 291 ---
\setcounter{footnote}{0}%
\opage{291}være; Konsten maa skabe Virkeligheden om og forvandle
Naturen.

Der er i slige Domme en egen Blanding af Sandt
og Falsk. Thi det er jo ganske sandt, at Konstværket
ikke kan være en ligefrem Kopi af den materielle Virkelighed,
den Virkelighed, som Enhver, der vil, kan
gaae hen at betragte med egne Øine; men det Usande
er, at Konstneren skulde have lagt Noget til af sit
Eget i den Forstand, at han til Forskjønnelse af Naturen
har anbragt Noget, som ikke hører med til Aabenbarelse
af Naturens Væsen, noget Ikke-Naturligt, altsaa
noget i Virkeligheden Unaturligt. Det er saa langt fra,
at den Natur, der har fundet sin Forklarelse i Konsten,
skulde være en forfalsket, med Virkeligheden stridende
Natur, at det tvertimod er Naturens inderste Væsen, den
sande Natur, der erholder sit Fænomen i det sande
Konstværk. Beskueren maa see Konstværket i det Lys,
hvori Konstneren har seet Naturen. Har han intet Øie
for Konstens Ide, saa seer han i Konstværket et vilkaarligt
Menneskeværk, en skuffende Efterligning, maaske
endog et Bedrag, og holder sig sandselig til den
sandselige Virkelighed. Hvad nytter det vel at vise
ham en Thorvaldsens Jason! Han seer en nøgen Mand
med en Hjelm paa Hovedet og et Faareskind paa Armen,
og spørger uvilkaarlig: har den virkelige Jason
seet saadan ud? Har Jason virkelig været til, og har
Thorvaldsen seet ham? Da man nu ikke kan give ham et
bekræftende Svar paa noget af disse Spørgsmaal, saa
erklærer han det Hele for et Fantasispil uden Virkelighed,
altsaa uden Sandhed.

Det, som ved slige Undersøgelser særlig bidrager
% --- p. 292 ---
\setcounter{footnote}{0}%
\opage{292}til at forvirre den sunde Sands, er en uklar Blanding
af de to Begreber: \emph{Virkelighed og Sandhed.} At
det Virkelige maa have sin Sandhed, og det Sande altid
paa en eller anden Maade svare til det Virkelige, er
uomtvisteligt; men derfor maa det ligefuldt indskærpes,
at Sandheden selv staaer uendelig høiere end den udvortes,
haandgribelige Virkelighed.

Dersom alt det Virkelige uden videre var Sandhed,
saa kunde der jo ikke gives nogen virkelig Vildfarelse,
nogen virkelig Forbrydelse eller virkelig Løgn: Vildfarelsens,
Løgnens, Forbrydelsens Virkelighed vilde da
kun være forskjellige Udtryk for virkelig Sandhed. Vistnok
taler man om en sand Forbrydelse, en sand Løgn;
men hvo seer ikke, at Tillægsordet »sand« i slige Forbindelser
netop stempler Forbrydelsen til Forbrydelse,
Løgnen til Løgn! Omvendt, hvis der ikke var Sandhed
uden i det, som er til paa udvortes haandgribelig Maade,
saa vilde der jo ikke være Sandhed hverken i Tanke
eller Fornuft, hverken i Sjæl eller Aand. Men at det
Indre, Usynlige ligesaa vel har sin Virkelighed og i
Virkeligheden sin Sandhed, som det Ydre, Haandgribelige,
ja en langt væsenligere og dybere Sandhed, vil
Ingen kunne nægte uden at fornægte al Sandhed.

At Sandheden i det væsenlige Indre er af høiere
Orden end Sandheden i det virkelige Ydre, er jo netop
det, som paa Ideens Trin giver Poesien sin Betydning
og guddommelige Ret. Poesiens stærkeste og dristigste
Brud med den ydre Virkelighed skeer i Eventyret. H.
C. Andersens Eventyr \emph{Skyggen} f. Ex. har mange interessante
Træk af det virkelige Liv og kunde næsten
faae Udseende af at være en virkelig Historie, dersom
% --- p. 293 ---
\setcounter{footnote}{0}%
\opage{293}den høitbegavede Forfatter ikke --- det gjør os ondt at
maatte sige ham det --- havde tilladt sig et saa høist
anstødeligt Brud paa Naturens uforanderlige Love; thi
siger en »fornuftig« Læser (og det Samme siger enhver
Dreng i en amerikansk Pogeskole): »det skal dog virkelig
Ingen bilde mig ind, at en Skygge saadan kan
liste sig bort fra sin Herre, krybe ind under Kagekonens
Skjørt, begynde at blive fed, medens Herren bliver
mager, og, hvad der er den værste Skandale og tillige
det Allerusandsynligste, drive det saavidt, at han tilsidst bliver forlovet, tænk --- med en Prindsesse af
Blodet.« Vi indrømme, at Forfatteren, vistnok af pur
Uvidenhed, hvad der kun tildels kan undskyldes med,
at han ikke er Fysiker af Faget, har overseet den meget
vigtige Omstændighed, at en født Skygge, selv om det
er Skyggen af »en lærd Mand«, aldrig, nei om den
end rækker sig og strækker sig nok saameget, aldrig
driver det til at blive et virkeligt Menneske, selv om den
blev et fedt Menneske. Kunde man bare tilgive Digteren
hans Uvidenhed i Fysik, saa vilde man ikke kunne nægte
ham Udmærket godt med Kryds og Slange i Psykologi;
thi de sjælelige Love, Indbildskhedens, Forfængelighedens,
Egoismens og Vindbeutleriets vel bekjendte Naturlove,
blive med den sande Genialitets ubedragelige
Sikkerhed paa det Allerpunktligste opfyldte i \emph{Skyggen.}

Vi behøve ikke her at fortabe os i vidtløftige Undersøgelser
angaaende Forholdet mellem Videnskab,
Konst og Poesi; thi Saameget vil af det Anførte tydelig
kunne sees: 1) at Spørgsmaalet om \emph{Virkeligheden}
altid maa underordne sig Spørgsmaalet om \emph{Sandheden;}
og 2) at Spørgsmaalet om Sandheden hver Gang
% --- p. 294 ---
\setcounter{footnote}{0}%
\opage{294}faaer en egen Betydning i Forhold til den Synskreds,
hvori Undersøgelsen falder. Sandheden kan i ingen Sfære,
ligesaa lidt i Poesien som i Videnskaben, opgive Virkeligheden;
men Sandheden staaer over den udvortes,
materielle Virkelighed. Det vigtigste Spørgsmaal er
ikke: hvormegen Virkelighed er der i det Sande; men
omvendt: hvormegen Sandhed er der i det Virkelige?
Naar det desaarsag atter og atter indskærpes, at hvad
der er Sandhed i Poesi og Konst, ikke er Sandhed i
Videnskab, og omvendt: saa er Meningen let at fatte.
Poesi og Konst kan ikke og bør ikke opfatte Virkeligheden
paa samme Maade som Videnskaben; den videnskabelige
Forsker kan ikke og bør ikke see Virkeligheden
med samme Øine som Digteren. Hvis en Zoolog
faldt paa at indføre Ænder og Gjæs talende, for ret at
fremhæve disse Dyrs karakteristiske Eiendommeligheder,
saa vilde man med Grund sige: det er en uvidenskabelig
Fremgangsmaade; der er ingen Alvor i en saadan
Behandling. Men fordi Zoologen synder imod Sandheden,
naar han lader Dyrene tale, hvo vil derfor paastaae,
at Eventyret: »\emph{den grimme Ælling}« ikke har
en dyb, til Virkeligheden svarende poetisk Sandhed.
At skille Skyggen fra sin Skyggegiver og lade den gaae
paa egen Haand er utilladeligt i Optiken; men der er
dog vel intet Menneske med Dannelse og sund Sands,
der beskylder Digteren for at krænke Sandheden, fordi
han af høiere Magtfuldkommenhed blæser Liv i et Skyggevæsen,
holder Læseren et Speil for Øiet og siger:
saaledes gaaer det til i Virkeligheden!

Det er ikke blot, som man siger, tilladt i Konst
og Poesi at forandre Noget ved det Virkelige og gjøre
% --- p. 295 ---
\setcounter{footnote}{0}%
\opage{295}Brud paa, hvad Naturlæren forstaaer“ ved Tingenes
Love; men det er en væsenlig Fordring, en uafviselig
Pligt saavel for Digteren som for Konstneren at behandle
Virkelighedsstoffet saaledes, at der kan være
Sandhed i Poesi og Konst. Ligesom Sandheden i alle
Kredse af Bevidsthedslivet har og maa have Virkeligheden
under sig, saaledes have Sandhedslovene i alle
Tilværelsens Kredse Virkelighedslovene under sig. Naar
blot Sandhedslovene opfyldes, saa falder Opfyldelsen af
Virkelighedslovene som af sig selv; sættes omvendt den
ydre Virkeligheds Love over Sandhedens Love, saa er
Konst og Poesi kun Løgn og Bedrag.

Overalt, hvor der spørges om Sandheden, spørges
der om en Dom over Virkeligheden, og da allerførst om
den Maalestok, hvorefter der skal dømmes. Hvad Videnskab,
Konst og Poesi angaaer, ere de Dannede i
det Hele dog saa temmelig vel orienterede; men, naar
Spørgsmaalet er om Forholdet mellem Sandhed og Virkelighed
i Religionen, kunde man fristes til at troe, at vi
endnu levede i Pavedømmets eller dog i Theologiens gyldne
Alder; thi Barbariet er afskrækkende. »Har Kristus
virkelig vandret paa Havet? har han virkelig opvakt
Lazarus? er han virkelig opfaren til Himmelen paa en
saa udvortes, synlig Maade, at jeg f. Ex., dersom jeg
havde været tilstede, kunde have seet det lige saa tydeligt
som Apostlene? Giv mig et ubedrageligt, historisk,
allerhelst naturhistorisk Bevis for, at det Altsammen
virkelig er gaaet til saaledes som det fortælles:
saa vil jeg troe det. Men dersom man ikke kan stille
mig fuld Sikkerhed for Virkeligheden af det Vidunder%
% --- p. 296 ---
\setcounter{footnote}{0}%
\opage{296}lige, saa maatte jeg jo være en Daare, om jeg antog
det for Andet end Sagn og Mythe«\footnote[1]{Hvorledes den udvortes Virkelighed kan i den hellige Historie stille sig afskrækkende op imod Aandsvirkeligheden, lære vi af Evangeliet om Kristi Vandring paa Havet. Af Frelserens Ord: \emph{Kom}! faaer Peder et Indtryk af Aandsvirkelighed. „Men der han saae det stærke Veir“, med andre Ord: da han fik et overvældende Indtryk af den sandselige Naturvirkelighed, „frygtede han“\dots{} Mtth. 14, 29—30.}.

De, der i vore Dage stille slige Fordringer til den
evangeliske Historie, røbe samme Mangel paa Dannelse
i Henseende til Religionen, som de, der i Sandhedens
Navn forlange, at Billedet af Jason bevislig skal ligne
den virkelige Jason, og Historien om Skyggen være en
virkelig Historie, røbe i Henseende til Poesi og Konst.
Løsreven fra Konstens Ide, faaer Virkeligheden ingen,
slet ingen Betydning i Konsten; løsreven fra Aabenbaringens
Aand, faaer Virkeligheden ingen, slet ingen Betydning
i Religionen.

Hermed siges ingenlunde, at det Virkelige skal opfattes
og bedømmes i Religionen paa samme eller dog
paa lignende Maade, som i Poesien og Konsten; langt
fra! her siges lige det Modsatte; thi her siges udtrykkelig,
at Virkeligheden skal opfattes og bedømmes paa
en anden Maade, efter en ganske anden Maalestok i
Religionen end i hvilkensomhelst anden af det aandelige
Livs forskjellige Sfærer, være sig i det praktiske
eller det theoretiske, være sig i Konst og Poesi eller i
Videnskab.

Hvad de empirisk forskende Videnskaber fastsætte
som strænge, uforanderlige Love, ere Love for den udvortes
givne Virkelighed. Saadanne Ting, saadanne

% --- p. 297 ---
\setcounter{footnote}{0}%
\opage{297}Love, og omvendt: ved denne Opfattelse bliver Videnskaben
staaende; hvorfra Tingene og Lovene have deres
Oprindelse, undersøger den ikke. Opkaste vi det Spørgsmaal,
om det ikke var muligt, at der kunde gives en
anden Tingenes Orden, en anden Natur end denne nærværende,
saa maa Videnskaben, naar den ret besinder
sig, svare: det er vel muligt, men saa maa der ogsaa
gives andre Love. Herved er Videnskabens Forhold til
Virkeligheden væsenlig karakteriseret. Religionen kan
meget vel indrømme Gyldigheden af denne Videnskabens
Grundsætning; men at den ikke kan indlade sig paa at
fremstille Virkeligheden i samme Belysning som Videnskaben,
ligger i Sagens Natur. Hvis Religionen betragtede
Virkeligheden fra samme Synspunkt som Videnskaben,
saa var den ikke Religion.

Sandhedslovene i Konst og Poesi ere Skjønhedslove.
Den hertil svarende Opfattelse og Behandling af Virkeligheden
behøver Religionen ikke at stille sig fjendtlig
imod; ja man kan vel endog sige, at den staaer i et
inderligere og fortroligere Forhold til Poesien end til den
videnskabelige Naturforskning. Ikke desto mindre lægger
Aabenbaringens Aand en ganske anden Vægt paa virkelige
Kjendsgjerninger og gjør anderledes strænge Virkelighedsfordringer
end den æsthetiske Ide\footnote[1]{Det æsthetiske Ideal tilhører alene Fantasien; det tilfredsstiller Kontemplationen; men stiller ingen Fordring til Villien; det kræver ingen Virkeliggjørelse, det har i Illusionen sin egen uvirkelige Virkelighed, sin Fuldendelse og Afslutning. Det religiøse Ideal er derimod i allerstrængeste Forstand et Virkelighedsideal, thi det er først og sidst et Villiesideal, og Villiesidealet kræver sin Virkeliggjørelse. Religionsfil. S. 205---6.}.

% --- p. 298 ---
\setcounter{footnote}{0}%
\opage{298}At Fysikere og Æsthetikere maaske kunne finde
Kategorien Aandsvirkelighed noget uklar, skal ikke undre
os; thi det er jo naturligt, at Enhver lever sig ind i
Sit og vænner sig til at see Virkeligheden med sit Øie;
men at de, der paatage sig at føre det store Ord om Religionens
videnskabelige Begrundelse, kaste Vrag paa Udtrykket
\emph{Aandsvirkelighed}, og --- ligesom til et Bevis
for, at deres egen Aand er uden Virkelighed, og deres
egen Virkelighed uden Aand --- forsøge at stemple dette
Udtryk til en Tvetydighed, forraader utvetydig, at de
i deres brændende Iver for Theologien have glemt den
første Troens Artikel, Artiklen om Skabelsen.

Troen paa Skabelsen forudsætter dog aabenbart, at
hele den synlige, sandselige Natur er Aandens Værk.
Men, naar den virkelige Natur oprindelig er Aandens
Værk, saa følger jo uimodsigelig heraf, at den aabenbarede
Religion vil have Alverdens Virkelighed betragtet
i Aandens Lys, at, med andre Ord, al Virkelighed maa
fra Begyndelsen af opfattes som Aandsvirkelighed. For
at undgaae Tvetydighed, er der da ogsaa paa sit Sted\footnote[1]{Religionsfilosofi. S. 475---76.}
udtrykkelig gjort opmærksom paa, at den religiøse Naturbetragtning
er i Principet absolut uensartet med den
videnskabelige. »Uden Hensyn til Videnskabens mulige
Indvendinger og længe før disse bleve til, har Aabenbaringen
udtalt sit Ord om Skabelsen, og «den religiøse
Bevidsthed har i Troens Inderlighed fundet begeistrede
Ord til Lovprisning af den Alt opholdende Magt, der
kan ryste Bjergene i deres Grundvolde, kan rulle Himlen
sammen som et Klæde, kan drage sin Aande tilbage, saa
% --- p. 299 ---
\setcounter{footnote}{0}%
\opage{299}at alt Levende døer\dots{} Det seer jo næsten ud, som
var den religiøse Naturanskuelse ligesaa sandselig objektiv
som den til Videnskaben svarende Iagttagelse; det seer
ud, som saa den Troende Gud selv i den ydre Natur,
i det mægtige Hav, hvor han i Vandene tømrer sine
Høisale, saae hans Øies Blink i Lynet, hørte hans Røst
i Tordenen, hans rullende Vogn i Uveirsskyer, naar
han ager paa Stormens Vinger. Dette er dog ikke saaledes,
det objektivt Naturlige er overalt kun Anledningsaarsag,
kun Middel og Gjennemgangspunkt for det i
Aanden Subjektive. «

Aabenbaringens Aand har Noget at meddele Adams
Slægt, Noget at oplyse Menneskene om, og dertil
bruger den Virkeligheden. Hvorledes den bruger det
udvortes Virkelige, i hvilke Kjendsgjerninger, Begivenheder,
Livsbilleder og Belysninger den anskueliggjør sin
Magt og Villie, sine Raadslutninger, Forjættelser og Love,
maa da ene og alene bedømmes efter det Maal, den har
sat sig, den Opgave, den i Religionen vil løse. Det er
i denne Henseende to Spørgsmaal, paa hvis Besvarelse
Alt kommer an: 1) \emph{hvorledes forholder Virkeligheden
af de i Bibelen fremstillede Begivenheder
sig til det Ord om Fald og Opreisning,
hvori Aanden aabenbarer den evige Sandhed?}
2) \emph{hvorvidt kan det stemme} med Aandens \emph{Ufeilbarhed
og Ordets Troværdighed, at der i den
hellige Historie anføres Træk af en Verdensbetragtning,
der vitterlig afviger fra og i enkelte
Stykker endog udtrykkelig strider imod
den i Nutiden gjældende Opfattelse af virke%
% --- p. 300 ---
\setcounter{footnote}{0}%
\opage{300}lige Naturforhold, af den i synlige Kjendsgjerninger
udtalte Tingenes Orden?}

Før der kan tænkes paa et ærligt og grundigt Forlig
mellem Religion og Kultur, er det nødvendigt at drage
disse Spørgsmaal frem og see, hvorvidt det efter Tidernes
Vilkaar er muligt at faae dem besvarede paa en for
begge Parter tilfredsstillende Maade.

Religionen staaer over Kulturen; havde Slægten kun
Valg mellem Eet af To: enten at opgive Religionen og
hylde Kulturen, eller holde fast ved den positive Religion
og bryde med Kulturen, da vilde de Alvorligere i
Samtiden vist ikke betænke sig paa at vælge det Sidste.
Vi holde os til dem, der strængt fordre, at det aabenbarede
Ord i alle Maader skal ske sin Ret; vi stille os
paa deres Side, som, hvis det var nødvendigt, hellere
vilde bryde Staven over den hele paa Naturvidenskaben
grundede Verdensbetragtning, end opgive en Tøddel af
det Troens Ord, der vitterlig hører til Frelse og Salighed.
Det er fra dette Synspunkt vi i denne Sammenhæng
ville yde vort Bidrag til Spørgsmaalenes Besvarelse.

Altsaa: \emph{hvorledes forholder Virkeligheden
af de i Bibelen fortalte Begivenheder sig til
Aabenbaringens Ord om Fald og Opreisning?}

For at begynde med Begyndelsen, henlede vi
Opmærksomheden paa Beretningerne i første Mosebog
om Skabelsen, Paradiset og Syndefaldet. At disse Beretninger
handle om noget Virkeligt og i allerstrængeste
Forstand angaae Virkeligheden, vil Ingen kunne nægte.
Vi sige: det er Aandsvirkelighed, Virkelighed for Aanden,
og forstaae dette saaledes, at den troende Opfattelse
aldeles ikke har med den udvortes fysiske, fra
% --- p. 301 ---
\setcounter{footnote}{0}%
\opage{301}Aanden løsrevne Virkelighed at gjøre, men at den Troende,
der i religiøs Selvforstaaelse vil forstaae Ordet, først og
sidst har med Aanden i Ordet at gjøre, har at betragte
de til Naturvirkeligheden hørende Træk som Elementer,
Midler og Betingelser, hvilke Aanden bruger for at anskueliggjøre
den Sandhed, den vil aabenbare.

For at godtgjøre Rigtigheden af denne Synsmaade,
skulle vi nu søge at vise, at to Mennesker kunne, hvad
Fortællingernes religiøse Indhold angaaer, troe det Samme,
det Selvsamme, med den samme Overbevisningens Styrke
og Inderlighed, og endda have en forskjellig Forestilling
om det i Begivenhederne udvortes Virkelige.

Lad os begynde med en enfoldig Troende. Han
opfatter da alt bogstaveligt; han følger Skabelseshistorien
Trin for Trin fra den første til den sidste Dag, og det
falder ham aldrig ind, at der kan være nogen anden
Forklaring af Verdens virkelige Tilblivelse. Han læser
Beskrivelsen af Edens Have, som man læser en Beskrivelse
af enhver anden virkelig Have, og gjør sig
ingen Betænkning over, at der maaske kunde reise sig
Tvivl angaaende Virkeligheden af de to overnaturlige
Træer, som her ere plantede i en forøvrigt naturlig
Jordbund. Han seer Eva foran Kundskabens Træ; han
hører Slangens fristende Ord; han bliver Vidne til Faldet
og føler i sin Samvittighed, at Jehovas Dom er ikke
blot stræng og forfærdelig, men ogsaa hellig og retfærdig;
kort sagt: den Troende opfatter i disse Begivenheder
det Udvortes i Enhed med det Indvortes, det Overnaturlige
i Enhed med det Naturlige, det Virkelige i Enhed
med det Aandelige, altsaa det Hele som Aandsvirkelighed,
Virkelighed for Aanden.

% --- p. 302 ---
\setcounter{footnote}{0}%
\opage{302}Stille vi os nu paa Dannelsens, Naturforskningens
og den ræsonnerende Kritiks Standpunkt, altsaa paa
det Bevidsthedstrin, hvor der kan være Tale om en
skarp, forstandig Adskillelse af det Aandelige og det
Virkelige som af det Overnaturlige og det Naturlige:
saa er det aldeles nødvendigt, at der fra Begyndelsen
af gjøres Rede for, hvad der for Religion og Tro er
det Afgjørende, om det er det udvortes Virkelige eller
det Indvortes, det Aandelige; thi paa Bestemmelsen af
dette Punkt kan Forholdet mellem Religion og Kultur
i det Hele siges at beroe.

Skulle Vidnesbyrdene om det udvortes Virkelige
sættes høiere end Aandens Vidnesbyrd, saa maa Naturvidenskaben
gjøres til øverste Dommer. »Det kan da
let bevises, at de i Skabelseshistorien anførte Kjendsgjerninger
ligesaa lidt lade sig bringe til Overensstemmelse
med, hvad der for Forskernes skarpe Blik har vist
sig at være Virkelighed, som Forestillingen om et Træ,
af hvis Frugter man skulde kunne spise sig til at blive
ligesaa klog som Gud selv, lader sig forlige med sund
Forstand. En Aand, der er saa ubekjendt med Naturlærens
og den sunde Fornufts allertarveligste Grundsætninger,
saa fuldstændig blottet for astronomiske, geologiske,
botaniske, zoologiske og anthropologiske Forkundskaber,
kan umulig være nogen virkelig Aand, endnu
mindre Guds egen Aand; hermed er Sagen afgjort.«

Men slige Angreb paa Religionen, med hvormegen
Skarpsindighed og naturvidenskabelig Kyndighed de end
udstyres, ere uden al Betydning, da de gaae ud fra en
grundfalsk Forudsætning. I Aabenbaring og Tro er det
ingenlunde det udvortes Virkelige, der staaer over Aanden,
% --- p. 303 ---
\setcounter{footnote}{0}%
\opage{303}men det er omvendt Aanden selv, der er ubetinget Herre
over al Virkelighed, og derfor bruger det Virkelige som
Middel for sit eget Øiemed. Grundspørgsmaalet er, i
hvilket Øiemed Aanden har udtalt sit Ord om Skabelse,
Paradis og Syndefald, hvad det er, den i dette sit Ord
vil undervise os om, og hvad Sandhed der er i denne
Undervisning.

Det er da for det Første Skabelsens Virkelighed,
hvorom det gjælder, men vel at mærke, opfattet som
Aandsvirkelighed. At den naturlige, virkelige Verden
er en Aabenbarelse af Guds almægtige, skabende Ord,
er Religionens og Troens oprindelige Forudsætning, en
Forudsætning, som den videnskabelige Kritik nok skal lade
være at rokke. For at gjøre denne Sandhed fattelig, anskuelig,
for de Enfoldige forstaaelig, blive de forskjellige
Skabelsestrin henførte til en Skabelse i sex Dage.

Ligesom Ordet om Skabelsen, saaledes har ogsaa
Ordet om Paradiset sin Virkelighed; men --- det er Aandsvirkelighed,
Virkelighed for Aanden. At Mennesket oprindelig
har Liv af Guds Liv, Aand af Guds Aand, og
altsaa er efter sin ideale Bestemmelse dannet til at bære
Guds Billede, erkjende Guds Bud og gjøre hans Villie,
dette er Anthropologiens øverste Sandhed, og den virkelige
Erfaring har nok som vist, at i samme Grad som
Bevidstheden herom fordunkles, i samme Grad synker
Mennesket, trods Forstandsdannelse og Kultur, ned i
dyrisk Afhængighed af selviske Naturmagter, og taber
Indtrykket af sin Slægts oprindelige Adel.

Og nu Ordet om Fristelse og Fald! Ethvert Menneske,
ogsaa den Helt, der er Gjenstand for Millioners
Beundring, begynder sin Løbebane med en dobbelt Fore%
% --- p. 304 ---
\setcounter{footnote}{0}%
\opage{304}stilling om Menneskelivet, en Forestilling om, at det i
Virkeligheden er, hvad det ikke burde være, og en
Forestilling om, at det i Virkeligheden burde være, hvad
det ikke er; denne Splid mellem Ideal og Virkelighed
kan den Enkelte kun see i Slægt og Historie, forsaavidt
han seer den i sit eget Indre. Helten kan nu gaae
to Veie. Troer han paa Ordet om det tabte Paradis,
saa erkjender han sig selv for medskyldig i Slægtens
Brøde, saa troer han paa Gudsbilledets Realitet i det
hellige Forbillede og anstrænger sig paa Opreisningens
og Befrielsens Betingelser af alle Kræfter for at arbeide
sig selv og Slægten op af den dybe Fornedrelse, den falske
Virkelighed, hvis Væsen er Trældom, Løgn og Bedrag,
og vinde et Skridt fremad mod den Aandens Virkelighed,
der fyldestgjør Livets Krav. Foragter han Ordet
om det tabte Paradis, saa vælger han en anden Vei.
I dunkel, mythisk Selvforgudelse begynder han da med
at gjøre det udvortes Virkelige med al den Forherligelse,
som derigjennem kan opnaaes, til Livets Maal.
Han kaster ikke ligefrem Foragt paa de ideale Forestillinger;
langt fra! Slægtens gamle Drømme, Paradisminderne
om Fred og Frihed, Ret og Sandhed, anvender
han til Vækkelsesmiddel for det godtroende Folk, til
ophøiede, uegennyttige Bevæggrunde, Trylleformularer,
hvormed man blænder Hoben. Han gaaer videre, han gjør
en kritisk Skilsmisse mellem det Virkelige og det Ideale.
Det Virkelige er, at den, der har Kræfterne, bruger dem
til sin Fordel; det Ideale er, at Verden vil bedrages;
saa smiler han til Verdens Aand og siger: ske din
Villie. En saadan Helteskikkelse maa nødvendigvis faae
to Ansigter, et officielt og et privat. Det officielle, der
% --- p. 305 ---
\setcounter{footnote}{0}%
\opage{305}vender udad imod Mængden, har ædle Træk, Præg af
Høihed og Majestæt; det er i ophøiet Ro, omstraalet
af Idealets Glands, et eftergjort, tilløiet Gudsbillede.
Men indenfor dette sidder et andet Ansigt, Personens
ærlige, virkelige Privatansigt, indadvendt med et Satyrgrin,
der spotter det udadvendte. Paa denne Maade
kommer Helten i Historien til at ligne Feen Alcine i
Eventyret: saalænge Trolddommen virker, og Elskerne
ere forblindede, seer hun ud som en deilig Kvinde; men
naar Trylleriet svinder, opdager man med Forfærdelse,
at det er en gammel, gusten, runken Hex. Men selv
denne Haan over Gudsbilledet og de hellige Paradisminder
er et talende Vidnesbyrd for, at Ordet om
Fristelse, Fald og oprindelig Skyld er Virkelighed,
Aandsvirkelighed, Virkelighed for Aanden.

At det nu ikke er Forestillingerne om det udvortes
Virkelige, om de sex Dages Skabelse, om Træerne i
Edens Have, om Slangens Udseende o. s. v., hvorpaa
det kommer an, sees tydelig deraf, at disse Forestillinger
ere som let henkastede Farvebilleder, der vel
kunne understøtte Anskuelsen, men ikke i og for sig
udøve den allerringeste Indflydelse paa Troen.

To Mennesker kunne --- vi gjentage det --- have den
samme Tro paa Ordets guddommelige Sandhed, og endda
være af høist forskjellig Mening angaaende det Udvortes
i Begivenhederne. Naar den Ene tager Alt bogstaveligt,
medens den Anden opfatter de ydre Træk sindbilledligt,
saa viser sig heri unægtelig en Meningsforskjel; men at
denne Meningsforskjel ikke angaaer Ordet som Aabenbaringsord,
kan sees, af to vitterlige Kjendsgjerninger.

For det Første er den, der tager Alt bogstaveligt,
% --- p. 306 ---
\setcounter{footnote}{0}%
\opage{306}ikke derved i Stand til at vinde noget nyt Udbytte for
Troeslivet. At troe paa Skabelsen er at troe, at denne
Verden virkelig er bleven til paa Almagtens Bud; men
Forestillingen om sex Dage, sex Uger, sex Millioner
Aar, gjør i religiøs Forstand hverken til eller fra. At
troe paa Paradiset og Syndefaldet er at troe paa Sandheden
af det Gudsbillede, hvori Mennesket oprindelig
er skabt, og i Selvmodsigelsens Smerte erkjende den
oprindelige Skyld, hvorved Idealet blev formørket. Men
Billedet af Kundskabstræet og den talende Slange kunne
kun ved en uhyre Misforstaaelse opstilles som Troesgjenstande.
At den Uskyldiges Fald forudsætter en
Fristelse, er en Troessandhed; men om den fristende
Slange har havt Ben og Menneskehoved, om den har
siddet paa Træet eller ligget ved Træets Rod, hvilken
Farve den har havt, om den har været blaagraa, grøn
eller spættet, er et Emne for digtende Fantasi, en
Meningsgjenstand, men ingen Gjenstand for Tro.

For det Andet vil Uenigheden mellen den, der tager
Alt bogstaveligt, og den, der opfatter Fortællingen sindbilledligt,
øieblikkelig vise sig som en Uenighed om
Fysik og Naturhistorie; men netop den Omstændighed,
at Undersøgelsesgrundene for og imod alle maae hentes
fra det udvortes Virkelige, ganske som ved Undersøgelserne
om Floden Phrat, om Bedellion og den Sten Onyx,
er jo et Vidnesbyrd om, at de Stridende have forladt
Aabenbaringens Grundlag og begivet sig ud paa Naturforskningens,
Troen uvedkommende, Omraade.

Men idet vi saaledes indrømme Aanden Ret til at
bruge Virkeligheden, som den vil, til Anskueliggjørelse
af sit Ords Sandhed, saa maa det for Alting ikke
% --- p. 307 ---
\setcounter{footnote}{0}%
\opage{307}glemmes, hvad det er for en Aand, hvorom Talen er.
Det er ikke en halvbevidst, uklar Menneske-Aand, som
her staaer afmægtig, famlende overfor den ydre Natur
og forvexler sine Drømmebilleder med virkelige Kjendsgjerninger;
nei, det er, ifølge Troens Forudsætning,
Guds almægtige Aand, den Livets og Sandhedens Aand,
der er Herre over al Virkelighed saavel paa Jorden som
i Himlen. Naar denne Aand i sit Ord bringer det
Virkelige ind under det Sindbilledlige, da holder den
Troende sig forvisset om, at det ikke skeer af Uvidenhed
eller Uklarhed, men i Sandhedens Navn for
Menneskets Skyld, fordi Mennesket, der trænger til
Gudserkjendelse længe før han paa Civilisationens Vei
er naaet saavidt, at han videnskabelig kan klare sig sin
Naturerkjendelse, kun i et saadant Sindbilledsprog er
i Stand til at fatte den guddommelige Tanke dybt og
rent. Herved er da tillige forudsat, at Aanden, forsaavidt
det tjener til Sandhedens Aabenbaring, ogsaa kan
bruge det Virkelige paa en anden Maade, at den i den
materielle Virkelighed kan indsætte en Virkelighed af
høiere Orden, at den, med andre Ord, kan aabenbare
sig ved Tegn og Undere.

Spørgsmaalet om Aabenbaringens Tegn og Under
er --- takket være Theologernes videnskabelige Maskepi
med den klare Atheisme --- i den nyere Tid blevet saa
forkvaklet, at den Dannede saare nødig indlader sig paa
videre Forhandlinger angaaende dette, efter hans Mening,
ligesaa ufrugtbare som forslidte Emne. Spørgsmaalet
indeholder nemlig to absolut uensartede Spørgsmaal,
det ene om Guds almægtige Villie, det andet om
Naturlovenes Uforanderlighed. Ved at sætte det sidste
% --- p. 308 ---
\setcounter{footnote}{0}%
\opage{308}Spørgsmaal øverst, beraabe sig paa Naturvidenskaben,
men fortie dens vitterlige Indrømmelse af vor Videns
absolute Grændse er det en let Sag at forvandle det
store, for Liv og Død afgjørende Grundspørgsmaal til
et theologisk tandløst Kællingespørgsmaal, som Ingen
længer gider høre paa.

De Hensyn og Indrømmelser, der i denne Henseende
skyldes Naturvidenskaben, ere, saavidt vi formaae at
trænge ind i Sagen, omhyggelig fremhævede i Religionsfilosofien\footnote[1]{Jvnf. Forfs. \emph{Religionsfilosofi}, S. 150---174; S. 476---477, o. fl. Stdr.}, og vi skulle her ikke videre opholde os ved
dette Punkt. Hvad vi i denne Sammenhæng maae lægge
Vægt paa, er, at Spørgsmaalet om Tegn og Undere
ubetinget maa henføres under Spørgsmaalet om den
guddommelige Villie. Er der virkelig en almægtig,
skabende Gud til og har denne almægtige, skabende Gud
en Villie? Paa dette enfoldige Spørgsmaals Besvarelse med
et enfoldigt Ja eller Nei, beroer den hele Undersøgelse.

At der i Filosofien er og maa være en Mængde
forskjellige Afskygninger i Opfattelsen og Bestemmelsen
af Gudsbegrebet, skulle vi være de Sidste til at nægte;
men fra Virkelighedens Synspunkt, paa Livets, Existensens
og den sædelige Bevidstheds Trin, staaer det urokkelig
fast, at Enhver, der nægter Gud en Villie, fornægter
Gud, thi en villieløs Gud er ingen Gud.

Det hele Mellemværende mellem Aabenbaringstroen
og den filosofiske Atheisme er let at opgjøre. »Vil Du,
det nittende Aarhundredes moderne Jeg, skyde Villien
ud af dit Gudsbegreb og leve dit Liv ud af Menneske%
% --- p. 309 ---
\setcounter{footnote}{0}%
\opage{309}slægtens metafysiske Princip, nu vel! det bliver din egen
Sag. Der er vist Ingen, der skal gjøre Forsøg paa
enten at tvinge Dig til Troen eller i nogen Henseende
paanøde Dig bibelske Forestillinger om det Overnaturlige.
Men saa bør Du ogsaa have en saa klar Bevidsthed
over dit Standpunkt, at Du ikke staaer og
bilder Dig selv og Verden ind, at det er en videnskabelig
Kritik af Troesprincipet, der har begrundet din
Gudsfornægtelse, medens det omvendt er din Gudsfornægtelse,
der danner Grundlaget for din opløsende Kritik
af den positive Religion.« Med denne korte Replik kan
en kyndig Religionslærer uden videre Besvær affærdige
Atheisten.

Hvad de enkelte Mirakler i den bibelske Historie
angaaer, maae Alle, som ere enige i at anerkjende den
første Troesartikel, ogsaa være enige i at anerkjende
disse Miraklers Mulighed; thi Troens ubetingede Forudsætning
er, at »for Gud skal ingen Ting være umuligt«.
Men fordi alle Undere antages for mulige, deraf
følger endnu ikke, at de Troende skulle være forpligtede
til at tillægge dem alle uden Forskjel samme udvortes
Virkelighed. Det er meget muligt, at et Tegn, en
Undergjerning, der efter vor Forestillingsmaade har Udseende
af at have dannet sig for Anskuelsen paa indvortes
Maade, er under Aandens Indvirkning bleven
opfattet og fremstillet, som om den hørte med i de
ydre Begivenheders Række. Hovedsagen er, at Aandsvirkeligheden
ligesaa vel omfatter det Udvortes som det
Indvortes, og at enhver ræsonnerende Adskillelse af det
Materielle og det Ideelle falder udenfor Troens Interesse.
Dette kan oplyses ved et eneste Exempel.

% --- p. 310 ---
\setcounter{footnote}{0}%
\opage{310}Har Kristus virkelig vandret paa Havet? Paa et
saadant Spørgsmaal, forsaavidt der udelukkende lægges
Vægt paa den ydre historiske Virkelighed, maa der
svares kort og bestemt: det er der intet Menneske paa
Jorden, der veed; thi om Sligt gives der ingen historisk
Viden. Verdenshistoriske Begivenheder, der alle
fremgaae af naturlige Aarsager, kan Videnskaben stadfæste
ved uforkastelige Vidnesbyrd; men om overnaturlige
Begivenheder kan der umulig gives nogen videnskabelig
afgjørende Vished, thi disse Begivenheder
maae ifølge Sagens Natur tage sig høist forskjelligt ud,
eftersom Øienvidnerne have seet dem i Troens eller i
Vantroens Lys.

Kristi Vandring paa Havet er ingen Virkelighed for
Videnskaben, men det er en Aandsvirkelighed, en Virkelighed
for Troen. »Thi,« siger den Troende, »hvorledes
en saadan Begivenhed skal bringes i Overensstemmelse
med Naturlovene, derom veed jeg Intet, derom
kan Intet vides; men dette erkjender jeg i Troen, at
for Gud ere alle Ting mulige, og dermed bortfalder al
Ræsonneren. Hvad der tiltaler mig i dette Livsbillede,
hvad der udelukkende fængsler min Opmærksomhed, er
det aandigt Ophøiede i Kristi Væsen, hans guddommelige
Opløftelse over de ydre Naturskranker, og saa
først og sidst den hjælperige Haand, hvormed han griber
den synkende Peder.« --- Dette religiøse Synspunkt er
i Principet absolut forskjelligt fra det theologiske.

Den theologiske Fortolker er Videnskabsmand; han
tænker som saa: »kunde jeg blot faae tilstrækkelig Vished
om, at det virkelig er skeet, som der fortælles; kunde jeg
ved Hjælp af en Smule Kritik, Hermeneutik og Exegetik
% --- p. 311 ---
\setcounter{footnote}{0}%
\opage{311}komme bag ved Aanden og faae fat i den nøgne Virkelighed,
saa vidste jeg da, hvad jeg skulde troe.« Vil
det lykkes Theologen? Ganske fortræffeligt! Han gaaer,
for ikke at blændes af Aandsvirkelighed, virkelig baade
bagved Aanden og udenom Aanden; og hvad finder han
saa? Et Sagn, en Mythe, en forvansket Virkelighed,
altsaa en Virkelighed, hvorpaa han ikke kan troe, at
sige, saalænge han staaer paa Kathedret; thi kommer
han engang paa Prædikestolen, ja saa skal Menigheden
vel faae at høre, at han er Mand for at troe paa den
nøgne --- og dog uforvanskede --- Virkelighed.

Hvad der i det Hele gjælder om Troeslærdommene,
gjælder ogsaa om de historiske Kjendsgjerninger, hvori
de anskueliggjøres, de ere Aabenbarelser af Aanden;
men just derfor tales der igjennem disse Kjendsgjerninger
ikke til det Sandselige i Mennesket, men til
Menneskets Aand. Den sandselige Virkelighed er Virkelighed
for Sandserne, det Forstandiges Virkelighed er
Virkelighed for Forstanden; men Aandsvirkelighed er
Virkelighed for Aanden. I den hellige Historie er Ordet
en guddommelig Kjendsgjerning, og Kjendsgjerningen et
guddommeligt Ord; den theologiske Mani bestaaer just
i at ville --- for Videnskabelighedens Skyld --- adskille,
hvad Gud har sammenføiet. Ved en saadan Adskillelse
maae Meningerne nødvendigvis faae Overvægt over
Troen, og usikkre Formodninger om det udvortes Virkelige
afgive den Maalestok, hvorefter Aanden skal lade
sig dømme. Men herom gjælder Pauli Ord: »det sandselige
Menneske fatter ikke de Ting, som høre Guds Aand
til, thi de ere ham en Daarlighed, og han kan ikke kjende
dem, thi de dømmes aandelig« (1 Korinth, 2, 14).

% --- p. 312 ---
\setcounter{footnote}{0}%
\opage{312}Og nu til vort næste Hovedpunkt:

2) \emph{Hvorledes kan det stemme med Aandens
Ufeilbarhed og ubetingede Sanddruhed
at fremstille Begivenheder i den hellige Historie
anderledes end de i Virkeligheden ere}
skete?

De Indvendinger, som heraf kunne udledes imod
det af os opstillede Begreb: Aandsvirkelighed, have ved
første Øiekast noget særdeles Slaaende.

»Naar« --- siger man --- »Gud i sit Ord virkelig
har villet aabenbare Menneskene, hvorledes det oprindelig
er gaaet til med Verdens Tilblivelse og Menneskelivets
Begyndelse, saa er det jo dog ganske utroligt,
at Aanden skulde have fremstillet Sagen anderledes end
den i Virkeligheden forholder sig, at den med Forsæt
--- thi af Uvidenhed kunde det da ikke være --- skulde
have anført Noget, som vitterlig er urigtigt, at, med
andre Ord, Sandhedens Aand ikke skulde have sagt den
fulde Sandhed.«

»Det er jo dog en underlig Selvmodsigelse at skulle
lade sig belære af Aanden, forsaavidt den taler om
Skabelsens Under, men i samme Aandedræt være
klogere end Aanden, forsaavidt den taler om Skabelsen
i sex Dage; at skulle bøie sig for Sandheden i det
aabenbarede Ord, forsaavidt det meddeler os en ellers
uopnaaelig Kundskab om Slægtens oprindelige Skyld,
men dog af Respekt for Geologi, Botanik og Zoologi
reservere sig den selvkloge Mening, at det i Virkeligheden
ikke forholder sig rigtig med Historien om Edens
Have, om Uskyldighedsstanden som en virkelig forhistorisk
Tilstand o. and. desl. Nei, Ærlighed varer
% --- p. 313 ---
\setcounter{footnote}{0}%
\opage{313}længst. Naar man ret vil være ærlig og tale rent ud,
saa er her kun Valg mellem Eet af To: enten erkjende
Aandens Ufeilbarhed i Et og Alt indtil Bogstav og
Tøddel eller forkaste Aabenbaringen og bygge paa Fornuften;
thi en Aand, der er ufeilbar i det Ene, men
ikke i det Andet, en Aand, som Skridt for Skridt hinker
imellem det Sande og det Usande, en saadan hinkende,
halvfeilende Aand kan ikke i noget Punkt gjøre Fordring
paa ubetinget Tro.«

Der er et gammelt indgroet Selvbedrag, der ligger
paa Bunden af disse foregivne Ærlighedsbetragtninger;
det er dette Selvbedrag, ikke Ærligheden, vi ville tillivs.

Det maa for hver ærligt Troende staae urokkelig
fast, at Aanden i sit Ord er absolut ufeilbar; det er jo
netop for at vække Erkjendelsen heraf, at vi saa eftertrykkelig
erklære os imod enhver ræsonnerende Adskillelse
mellem Tanken og Indklædningen. Men Indklædningen
er i Aandens og Sandhedens Navn lige ufeilbar,
enten Klædningsstykkerne, om vi saa maae sige, ere
hentede fra den lavere eller den høiere, den sandselige
eller den oversandselige, den jordiske eller den himmelske
Virkelighed, enten Klædningen er vævet af immaterielle
eller materielle Traade.

Det fromme Selvbedrag, hvortil her sigtes, bestaaer
i, at Mennesket gjør det til Hovedsag, hvad Aanden i
sin Visdom har nedsat til Bisag, at den Troende, uden
selv at vide deraf, sætter det Afgjørende i Noget, hvori
det ikke skal sættes, og selvraadig foreskriver Aanden
en Ufeilbarhed, som den ikke vil vedkjende sig. De,
der have lidt Kjendskab til Kirkens og Theologiens
% --- p. 314 ---
\setcounter{footnote}{0}%
\opage{314}Historie, have Exempler nok paa Vildfarelser og Misgreb
i denne Retning\footnote[1]{Jvnfr. \emph{Om „den gode Villie“} som \emph{Magt i Videnskaben}. Kbhvn. 1867. S. 40—75.}.

For at sikkre sig det gamle Testamente som fuldkommen
paalideligt Aabenbaringsdokument, gik protestantiske
Theologer i det 17de Aarhundrede ud fra den
Forudsætning, at endog Vokalpunkterne skyldtes Inspirationen.
De holdt fast ved den Paastand, at en ufeilbar
Aabenbaring krævede en ufeilbar Text; men hvor usikker
Texten er, naar man kun har Konsonanter og vilkaarlig
kan tilføie Vokalerne, er indlysende af sig selv\footnote[2]{Saaledes vil man af f. Ex. Stedet i 1ste Mosb. (Kap. 23, 4), hvor Mødet mellem Esau og Jakob omtales, faae ud, efter som man forandrer Vokaltegnene, baade at Esau kyssede Jacob og at han bed ham.}. De
lærde Buxdorfer i Basel, Fader og Søn, strede tappert
for Aandens Ufeilbarhed, idet de strede for Vokalernes
Ælde; men sproglig-historiske Kjendsgjerninger lade sig
ikke bortræsonnere. Stridens Udgang viste, at de orthodoxe
Ivrere ikke havde Aanden paa deres Side, at de
i det omstridte Punkt havde dannet sig et vilkaarligt
og falsk Begreb om Ufeilbarheden.

Paa lignende Maade forholder det sig med den
theologiske Strid om den Renhed og Klassicitet, man
har villet tilkjende det græske Sprog i det nye Testamente.
»Det er Aanden, som paa overnaturlig Maade
har indgivet Apostlene Sproget; det nye Testamentes
Græsk maa altsaa være klassisk Græsk. Er der Nogen,
der trods Plato, Herodot og Xenofon kjender det græske
Sprog fra Grunden af, saa er det nok den Hellig Aand.«

% --- p. 315 ---
\setcounter{footnote}{0}%
\opage{315}At drage denne Paastand i Tvivl var, ifølge Puristernes
Mening, at drage Aandens Ufeilbarhed i Tvivl. Naar
det nu ved sindig, grundig, alsidig Sprogforskning ikke
destomindre er blevet godtgjort, at Græciteten i det nye
Testamente ikke er klassisk, men hellenistisk med kjendelige
Indblandinger af hebraiserende Talemaader, hvad
beviser saa det? Det beviser, at Puristerne have forfeilet
Synspunktet for den rette Opfattelse af Sproggaven og
søgt Fuldkommenheden der, hvor den ikke skal søges.

Det samme Selvbedrag gjentager sig ved Bedømmelsen
af de historiske Beretninger. Endnu i sidste
Halvdel af det 18de Aarhundrede udtaler Theologen
Chr. W. F. Walch\footnote[1]{Krit. Nachricht v. d. Quellen der Kirchengeschichte. Leipz. 1770. S. 70.} sig om de hellige Skribenters
Ufeilbarhed paa følgende Maade: »de have ikke alene
ikke feilet, men de have« --- nemlig paa Grund af Inspirationen
--- »end ikke kunnet feile, og deres Fortællinger
ere desaarsag ikke blot ligesaa visse som enhver
anden historisk Sandhed, men endogsaa ubedragelig
visse.« Dersom Historikeren nu herved alene tænkte
paa Ordet som Aabenbaringsord, d. v. s. paa Sandhedsindholdet
i dets uopløselige Enhed med Form og
Indklædning, altsaa netop paa, hvad vi kalde Aandsvirkelighed,
Virkelighed for Aanden, saa vilde den videnskabelige
Kritik, hvormeget den end havde bidraget til
at forandre Meningerne og Formodningerne om det Udvortes
og Endelige i Fortællingernes Sammensætning,
ikke have været istand til at forandre en Tøddel i Ordet
selv\footnote[2]{Om Aabenbaring og hellig Historie. Jvnfr. Religionsfil. S. 24—28.}. Men saaledes forholder det sig ikke. Istedenfor
% --- p. 316 ---
\setcounter{footnote}{0}%
\opage{316}at anlægge Aandens Maalestok paa det Virkelige, gaaer
han bagvendt tilværks og anlægger Virkelighedens, den
udvortes endelige Virkeligheds Maalestok paa Aanden.
Han trodser paa den Ufeilbarhed, han selv paadutter
Aanden, og aner ikke, at han netop derved giver Ordet
til Pris for opløsende Kritik\footnote[1]{Et Par Exempler. I anden Samuels Bog (XXIV) berettes Følgende: „Og Herrens Vrede blev atter optændt mod Israel, og han tilskyndede David imod dem, idet han sagde: Gaa, tæl Israel og Juda!“ Parallelstedet til anden Samuel (XXIV) er, som bekjendt, første Krønikerne (XXI), der begynder saaledes: „Satan stod Israel imod og tilskyndede David til at tælle Israel.“ Opfattet som Fortælling, staaer den ene af disse Beretninger i bogstavelig Modsigelse med den anden. Men i Aabenbaringsordet selv er der ingen Modsigelse; thi Tilskyndelsen, der, ifølge den første Beretning, udgaaer fra Herrens Vrede, henføres i den anden Beretning til den onde Villie, der netop i Vredesforholdet stiller sig mellem Jehova og David. Saaledes er det i Joh. 13, 26—27 jo baade Kristus og Satan, der tilskynde Judas til Forræderiet; men hvo vil heri see en religiøs-psykologisk Modsigelse.}.

Henvende vi nu igjen Opmærksomheden paa hvad
der i det Foregaaende er anført om Skabelse, Paradis
og Syndefald: da ville alle til Selvbedraget hørende
Misforstaaelser samle sig i fuld Belysning. Den enfoldig
Troende, der tager Alt efter Bogstaven, fordi han ikke
gjør sig nogen Forestilling om, at nogen anden Opfattelse
var tænkelig, har ingen Skyld i Misforstaaelserne.
Misforstaaelserne, det maa udtrykkelig fastholdes, begynde
først ved det Vendepunkt, hvor den ræsonnerende
Forstand søger Fodfæste i det Virkelige og vil ved at
stille det Naturlige udvortes sammen mod det Overnaturlige
kontrollere Aanden.

Guds Aand stilles da udvortes imod Menneskets
Aand; Gud forestilles at tale til Mennesket paa samme
% --- p. 317 ---
\setcounter{footnote}{0}%
\opage{317}sandselige Maade, som naar det ene Menneske taler til
det andet. Aanden dikterer de hellige Mænd, hvad de
skulle skrive; er der Feil i Diktatet, saa maa Aanden
naturligvis bære Skylden: det er det stærkeste Vrængebillede
af Aandsvirkelighed, som vel lader sig tænke.
Grundvildfarelsen er, at det Overnaturlige behandles,
som om det var naturligt.

I Aabenbaringen er Forholdet mellem Gud og Menneskets
Aand et Aandsforhold; Aanden taler paa indvortes,
ikke paa udvortes Maade; selv naar Indtrykket
er saa stærkt, at Ordet lyder for Troens Øre, udspringer
det Udvortes dog af det Indvortes. Derfor taler Aanden
et menneskeligt Sprog; den udlægger sit Ord netop i
det Sprog, som de Mennesker forstaae, hvortil den
taler; derfor talede den hebraisk til Jøderne, hellenistisk
Græsk til Datidens Grækere, og vil endnu den Dag
idag oprindelig tale Dansk til Danske, Tydsk til Tydske
og Fransk til Franske.

Opfattes Forholdet mellem Guds og Menneskets
Aand som et virkeligt Aandsforhold, saa vil man ogsaa
kunne forstaae, at Aabenbaringens Ord ikke blot iklæder
sig det menneskelige Sprog, men i det Hele retter
sig efter Menneskets Forestillingskreds, betjener sig af
dets egen Verdensbetragtning, dets naturlige Syn paa
Himmel og Jord, for ved disse Tilknytningsmidler at
anskueliggjøre den Sandhed, kvorpaa det alene kommer   % PRINTED AS IS (Bind 3, p. 317): „kvorpaa“ for hvorpaa (zoom z4)
an, Troessandheden til aandig Vækkelse, til Fred og
Frelse.

Skulde det fra dette Synspunkt nu dog ikke være
indlysende, hvilket Selvbedrag der ligger i den Fordring,
at Aanden i sit Ord om Skabelse, Paradis og Synde%
% --- p. 318 ---
\setcounter{footnote}{0}%
\opage{318}fald skulde have lagt de astronomiske, geologiske og
naturhistoriske Forestillinger til Grund, som vise sig at
have Gyldighed i den nyere Videnskab! Vi maae jo
skamme os ved en saa egensindig, uforstandig Afgjørelse
som denne: enten er det Aanden, der har taget feil, og
da er Ordet ikke troværdigt, eller Aanden selv maa
være ufeilbar ogsaa i Beskrivelser af det Fysiske, og
da maae vi see at indprænte den kjære Ungdom al
mulig Mistro til Videnskaben.

Der staaer skrevet: værer enfoldige som Duer;
men der staaer ingensteds skrevet: værer dumme som
Flyndere. Enhver, der i vor Tid vil bruge sin sunde
Forstand og ikke i aandeligt Hovmod og taabelig Trods
kaste Vrag paa menneskelig Viden, vil, hvad Verdensbygningen
angaaer, let overbevise sig om, at Forargelsen
paa det kopernikanske System beroer paa en
selvklog Indbildning. Lad det i astronomisk Forstand
være tusinde Gange sandt, at det er Jorden, der gaaer
omkring Solen, det er da ligesaa sandt, at Aabenbaringens
Aand umulig har kunnet begynde med at antyde
dette, uden at gjøre sit Ord aldeles uforstaaeligt for
Menneskene.

Man lægge vel Mærke til, at Oplysning om en saa
indviklet Sag ikke er given med en enkelt Bemærkning.
Alene det at skulle forflytte sit Standpunkt fra Jorden
til Solen og bestemme Stjernernes heliocentriske Steder,
hvilken Anstrængelse koster det ikke Begynderen. Der
er iblandt de saakaldte Dannede Tusinde og Tusinde,
der »troe paa« det kopernikanske System, men som vilde
komme i den pinligste Forlegenhed, hvis de skulde gjøre
Rede for Beviserne. Og saa skulde Aanden endda
% --- p. 319 ---
\setcounter{footnote}{0}%
\opage{319}have bevist sin astronomiske Ufeilbarhed ved at gjøre
det begribeligt for Abraham og Moses, som endnu neppe
er bleven begribeligt for tusinde Parten af dem, der
drive med Dannelsens Strøm.

I sit Ord er Aanden ufeilbar; til Sandhedens
Aabenbaring i og ved Ordet tager den ubetinget al
Virkelighed, baade den udvortes og den indvortes, i
Tjeneste; derfor kan Troens Øie umulig see det Virkelige
for sig i et blot naturligt Lys.

Udenfor Aabenbaring og Tro er der intet Paradis
og har aldrig været; men i Aabenbaringens Ord for
Troens Blik har Paradiset fuld Virkelighed. Den Troende
seer Alt i det hellige Lys, og Alt er virkeligt: Edens
Have, Livets Træ og Gihons Flod; han seer baade
Slangen og Kvinden. Men i Videnskabens Lys er her
Intet at see, slet Intet.

Fra et geografisk Stade betragter man Ørken Sinai:
Granitbjerge med vilde, vidunderlig takkede Spidser rage
frem over vandrige, frugtbare Dale; man seer Horeb,
Dschebel Muse, den egenlige Sinaispids; det er Altsammen
virkeligt: her er Virkelighed nok for Geografien.
Men for den Slags Virkelighed er Troen blind; den er
i Geografens Øine som Uglen, der ikke kan see om
Dagen. Det kommer af, at Troens Øie er anderledes
organiseret og trænger til et andet Daglys end Naturbetragterens.
Først, naar Aandens Lys bestraaler Sinai
Bjerg, bliver Troen seende; saa seer den virkelig Bjergets
Top; den seer Jehova og Moses: i Tordenskyen er
Jehovas Finger og Tavlerne, hvorpaa han skriver!
Men, mærkværdigt nok, nu kan Naturforskeren Intet see;
nu er det ham, der er blind som en Ugle.

% --- p. 320 ---
\setcounter{footnote}{0}%
\opage{320}At Genezareths Sø er et luende Speil, naar Solen
gaaer op, og Straalerne gløde paa Thabors Fjeld, er en
Landskabsherlighed, som endnu aldrig er seet med Troens
Øie. Det betyder ingenlunde, at den Troende skulde være
lyssky eller mangle Sands for Naturskjønheder; men det
betyder, at Troen som Tro maa have sit eget Lys.
Naar dette Lys er tændt, seer den Virkeligheden overalt,
den naturlige saavel som den overnaturlige; den seer
Bjerget og den Forklarede paa Bjerget; den seer i
Nattens Dæmring Genezareth ved Søen og Bølgerne og
Skibet og Frelseren vandrende paa Søen. Det er Aandsvirkelighed,
Virkelighed for Aanden.
\signature{R. \emph{Nielsen}.}

\clearpage

% --- front matter: the title page of Bind 5 (PDF 1 of 1871-2-for-idee-og-virkelighed.pdf), read at the image (reader K, m1);
% short wavy rules are set as plain rules; the library stamps are not transcribed.
\begin{titlepage}
\begin{center}
\vspace*{3em}
{\Huge For Ide og Virkelighed.}\\[1em]
\rule{8em}{0.4pt}\\[1em]
{\large Et Tidsskrift,}\\[0.8em]
{\small udgivet}\\[0.4em]
{\footnotesize af}\\[0.6em]
{\large R. Nielsen,}\\[0.4em]
{\large Bj. Bjørnson og Rud. Schmidt.}\\[1em]
\rule{8em}{0.4pt}\\[0.6em]
{\large Aargang 1871. Andet Bind.}\\[0.6em]
\rule{8em}{0.4pt}\\[4em]
Kjøbenhavn.\\[0.4em]
Otto B. Wroblewskys Forlag.\\[1.5em]
{\small G. S. Wibes Bogtrykkeri.}
\end{center}
\end{titlepage}

% ===== PART II: Bind 5, pp. 1-38 =====
\essayhead{Videre om Aandsvirkelighed,}{foranlediget ved Pastor V. Birkedals og Pastor C.\\ Møllers Indlæg.}
% --- p. 1 ---
\setcounter{footnote}{0}%
\opage{1}Enhver, som med samlet Sind og uhildet Blik vil gjennemgaae
Beretningerne i det Gamle og det Nye Testamente
og søge at gjøre sig de overnaturlige Begivenheder
nærværende, maa vistnok, dersom han da ikke mangler alle
Dannelsens Forudsætninger og Selvtænkningens Betingelser,
faae underlige Tanker. Hvor har dog ikke den hele
virkelige Tingenes Orden seet anderledes ud i Profeternes,
Kristi og Apostlenes Dage, end den seer ud i vore!
Det er jo, som talede den hellige Skrift om en anden
Himmel og en anden Jord, en anden Verden og en anden
Gud, kort: om en ganske anden, ikke blot indvortes, men
ogsaa udvortes Virkelighed end den, vi virkelig kjende.
Den speidende, prøvende og paa Virkelighedsiagttagelser
ellers saa sikkre Nutid, om den end anstrenger sin Fornufts
Kræfter til det Yderste, er ikke i Stand til at bringe
sit Begreb om Virkelighed og virkelige Kjendsgjerninger
i Samklang med, hvad Bibelen kalder det Virkelige. At
tye til Videnskaben om Hjælp og Bistand har ganske vist
sin Betydning, naar Opgaven er at skille sig af med den
% --- p. 2 ---
\setcounter{footnote}{0}%
\opage{2}bibelske Virkelighed; thi hvorledes Fremstillinger af det
Overnaturlige lade sig opløse i Sagn og Myther, derom
veed Videnskaben fortræffelig Besked, og Ingen, hvem
de mirakuløse Begivenheder falde til Besvær, skal fortryde
at have søgt dens Veiledning. Men at ville --- i
den Indbildning, at Troen er ensartet med Viden --- plage
Videnskaben, at den skal gaae i Kaution for Tilforladeligheden
af det Mirakelvæsen, der er saa anstødeligt for
»alle Dannede«, kan naturligvis kun falde dem ind, som
endnu have en theologisk Skrue i Hovedet.

Skulle de naturlige Begivenheder i den profane Historie
og de overnaturlige i den hellige Historie virkelig
sees i samme Lys og undergives samme Kritik, saa er
Kristendommen tabt for Civilisationen: Bibelens Troværdighed
staaer ikke til at redde. For at komme ud af
det Virvar, hvori en objektiv ubarmhjertig Kritik i Videnskabens
Navn har indfiltret den ulyksalige Theologi, saa
den hverken veed ud eller ind, hverken tør være kritisk
eller ukritisk, hverken opgive Miraklet eller holde paa
Miraklet, gives der kun een Udvei: ubetinget at erkjende
det religiøse Princips Selvstændighed og erklære den
hellige Historie for uensartet, \emph{absolut} uensartet, med
den profane.

Er det først tilfulde indseet, at de bibelske Beretninger
maae, ifølge Sagens Natur, være af en anden Oprindelse,
grunde sig paa en anden Virkelighed og falde
ind under en anden Synskreds end Beretninger, der høre
hjemme i en almindelig Verdenshistorie, saa følger jo
uimodsigelig heraf, at de bibelske Beretninger maae sees
i deres eget Lys og bedømmes fra deres eget Standpunkt.
Som Virkeligheden er, maa Kritiken være. For da at
% --- p. 3 ---
\setcounter{footnote}{0}%
\opage{3}give hin høiere Virkelighed et betegnende Navn har jeg
kaldt den Aandsvirkelighed, Virkelighed for Aanden, idet
jeg ved det første Udtryk særlig har sigtet paa Kjendsgjerningernes
guddommelige Ophav, ved det sidste derimod
paa den Bevidsthed, der skal forvisse sig om
Kjendsgjerningernes Sandhed. Jeg skal nu søge at bringe
den herved antydede Hovedtanke saa nær til Øiet som
muligt.

\begin{center}I\end{center}

\emph{Aandsvirkelighed afgiver den første nødvendige
Forudsætning for Bedømmelsen af de
hellige Kjendsgjerninger, naar den religiøse
Bevidsthed har naaet det Dannelsestrin}, da
\emph{Videnskabens fornuftige Indvendinger kunne}
gjøre Indtryk paa \emph{den, saa at den ikke blot udefter},
\emph{men ogsaa indefter maa befinde sig i
Strid med den videnskabelige Kritik}.

»Virkelighed for Aanden«, siger Pastor Birkedal, »er
ikke den første Virkelighed«. Virkeligheden maa »først
bevise sig for en anden Modtagelsesevne (Organ), før den
kan bevise sig for Aanden, hvis Kristendommen virkelig
skal blive vor Frelse, vort Liv. Ganske vist maa det
skee ved Aanden, fordi Kristendommen er et Aands-Rige,
hvori Intet kommer og skænkes \emph{uden} Aanden, og
forsaavidt kan man kalde Kristendommens Virkelighed
»Aandsvirkelighed». Men Forklaringen af dette
Udtryk: »Virkelighed for Aanden« synes mig desuagtet
ikke at svare til det »virkelige« Forhold, naar der skal
begyndes dermed«. Ja, naar der skal begyndes dermed!
Heri har Pastor B., efter min Opfattelse, fuldkommen Ret.
% --- p. 4 ---
\setcounter{footnote}{0}%
\opage{4}Men Sagen er, at der hverken saaledes er begyndt eller
skal begyndes dermed. Om Troeslivets Oprindelse ere
vi, hvad Pastor C. Møller ogsaa gjør opmærksom paa, i
alt Væsenligt vistnok aldeles enige. Her er i denne
Sammenhæng ikke Spørgsmaal om, hvorledes Troen kommer
ind i et Menneskes Hjerte, eller hvorledes et Menneske
kommer til Troen; men Spørgsmaalet er, hvorledes
et Menneske, som af Hjertet troer paa Guds aabenbarede
Ord og med denne Tro har levet sig ind i Bibelens hele
Forestillingskreds og fundet Trøst og Glæde i Forvisningen
om, at de hellige, overnaturlige Kjendsgjerninger ikke
ere svævende Sagn og Billeder, men fuld tilforladelig
Virkelighed, hvorledes, siger jeg, et saadant Menneske vil
være i Stand til at fastholde denne Virkelighed lige over
for en Kritik, som med vægtige videnskabelige Grunde
indskærper og under mange Vendinger gjentager hin fra
den negative Viden vel bekjendte Sætning: Det \emph{Guddommelige
kan ikke være skeet saaledes; det
saaledes Skete kan ikke være guddommeligt}.

Det er ikke Agtelse for Kjendsgjerningerne, ikke
Kjærlighed til Bibelordet eller et grundigt Studium af
den hellige Skrift, der gjør Skrifttheologien aandløs; det
Aandløse, vil sige: det Theologiske, ligger netop i den
ulykkelige Misforstaaelse, hvormed man begynder. Uden
at gjøre sig Rede for, at det er en ganske egen Art af
Virkelighed, man her har for sig, gaaer man stiltiende
ud fra den uklare Forudsætning, at Virkelighed er Virkelighed,
Begivenhed Begivenhed; --- Materialist er jo Materialist,
som det hedder i Vaudevillen. Skrifttheologien
stiller, med andre Ord, Virkeligheden af det Overnaturlige,
vel at mærke: hvad Virkelighedskriteriet angaaer,
% --- p. 5 ---
\setcounter{footnote}{0}%
\opage{5}paa lige Trin med Virkeligheden af det Naturlige, og
hvad bliver Følgen? De samme historiske Beviser, der
anvendes til at forsvare de naturlige Begivenheder, anvendes
nu ogsaa til at forsvare de overnaturlige, den samme
Sandsynlighedsberegning, der anlægges paa hine, skal nu
ogsaa passe paa disse. Saa møder den bibelske Theolog
med, hvad han kalder »et troende Sind«; det vil sige: han
er villig til at indrømme, at der er skeet noget Underfuldt,
men hvori det Underfulde egenlig bestaaer, skal
først udfindes ved fortsat videnskabelig Granskning; han
er villig til at antage Mirakler, naar de ikke ere ham altfor
urimelige og usandsynlige. Men desværre for ham, der
er i den hele Bibel ikke et eneste Mirakel, nei ikke et
eneste, som jo ikke fra et videnskabeligt Standpunkt, nøie og
nær beseet, er urimeligt, usandsynligt og frastødende.
Den skrifttheologiske Kritik tager en andægtig Mine paa,
naar den skal forsvare Opstandelsens og Himmelfartens
store Kjendsgjerninger, hvorimod den ganske bevarer sin
videnskabelige Holdning over for Miraklet med Gadarenernes
Svin; den tænker som saa: med Svin og Djævle
behøver man ikke at gjøre Komplimenter. Men hvor er
nu Troen, og hvor er Kritiken?

At Opstandelsens Under har en ganske anden gjennemgaaende
Betydning for det religiøse Liv end Begivenheden
med den Besatte i Gadara, kan det kun falde en
Afsindig ind at ville nægte; men her handles nu ikke
om et Mirakels større eller mindre religiøse Betydning;
her handles om Miraklets Virkelighed og den mirakuløse
Virkeligheds større eller mindre Sandsynlighed. Naar
dette Synspunkt fastholdes, og det maa fastholdes strengt,
dersom Kritiken ellers skal have Noget at betyde, saa
% --- p. 6 ---
\setcounter{footnote}{0}%
\opage{6}behøves der ikke megen Skarpsindighed for at indsee, at
alle Mirakler ere lige sandsynlige d. v. s. lige usandsynlige,
de ligne hinanden ubetinget i Eet: de ere alle Mirakler.
Opstandelsens Mirakel kommer ikke den ligefremme
historiske Sandsynlighed et Haarsbred nærmere
end Miraklet med Dæmonerne i Svinehjorden.

Jo mere Bibeltheologen med Videns Fakkel vil lyse
op i den dunkle Grund, der bærer de hellige Kjendsgjerninger,
desto mørkere, tættere Taager opstige der fra
Dybet, og desto uhyggeligere maa han blive tilmode.
Der er i denne Omgivelse overalt et Noget, der er hans
Bevidsthed imod, et Uforklarligt, der trykker ham for
Brystet og tynger paa hans videnskabelige Samvittighed
--- hvis han da virkelig har en videnskabelig Samvittighed.
Herfra den grændseløse Tvivlraadighed overfor det
mest Afgjørende; herfra den endeløse Vaklen mellem Antagelse
og Ikke-Antagelse, den afsløvende Halvvering og
intetsigende Forklaring, at der ganske vist er skeet Noget,
noget Overordenligt; men hvad det er, der er skeet, det
veed man ikke. Saa bliver Aabenbaringens Egn et Taagernes
Land, dens Indhold et bølgende Element, hvori
man gaaer tilbunds i Videnskabens Dykkerklokke.

Theologien maa bort. Her maa vælges et andet
Udgangspunkt; her maa begyndes ud fra en anden Grund.
Det maa en Gang for alle være afgjort, at den, der træder
ind i Aabenbaringens Lysregion og føler sig tiltalt af
de overnaturlige Fænomener, har overskredet Videns
Grændse og sagt Farvel til de tvingende, almennødvendige
Fornuftgrunde. Dermed har den Troende ingenlunde
enten opgivet Tankelovene eller udslukt Fornuftens Lys;
men Eet har han opgivet, den falske Begrundelse, Be%
% --- p. 7 ---
\setcounter{footnote}{0}%
\opage{7}grundelsen i objektiv Viden. Han føler og forstaaer, at
Religionen umulig kunde være Religion, dersom den ikke
havde sin egen Oprindelse, sit eget selvgyldige Princip;
han føler og forstaaer, at Kjendsgjerninger, der fremgaae
af et overnaturligt Princip (den hellige, almægtige Villie),
maae have et tilsvarende overnaturligt Virkelighedspræg;
han føler sig forvisset om, at den hele, i sig sluttede
Erkjendelseskreds, hvormed hans Sjæl er fortrolig, og
hvori hans bedre Væsen føler sig hjemme, tilhører en
Virkelighed, der lige saa lidt kan stadfæstes som undergraves
af Videnskaben. Med dette Fodfæste paa Virkelighedens
overnaturlige Enemærker er han som en Troende,
og dog som en Vidende, forberedt til at tage imod den
videnskabelige Kritiks Indvendinger. Disse Indvendinger
ville lige saa lidt være i Stand til at rokke hans Tro paa
Aabenbaringen som at svække hans Agtelse for Videnskaben.
Det er saa langt fra, at han i Religionens Navn
vil forbittres paa den objektive Kritiks Angreb, at han
meget mere vil indsee, at den med sine Forudsætninger,
sit Grundlag og sit Virkelighedsmaal umulig paa noget
Punkt kan tilegne sig det Overnaturlige, at den i saa
Henseende er lovlig undskyldt. Det vil i det Hele som
i det Enkelte efterhaanden blive ham ret indlysende,
hvorledes det hele Væv af kritiske Vanskeligheder væsenlig
skyldes theologiske Vævere, som have gjort Begyndelsen
med at sammenvæve to Slags Virkelighed, fysisk
Virkelighed og Aandsvirkelighed.

\begin{center}II\end{center}

\emph{Aandsvirkelighed udtrykker Virkelighedsbegrebet
i dets hele Omfang og omfatter saa%
% --- p. 8 ---
\setcounter{footnote}{0}%
\opage{8}ledes baade det indvortes og det udvortes
Virkelige, baade den legemlige og den ulegemlige
Virkelighed, men altid med særligt Eftertryk
paa Aanden selv, den almægtige Villie},
\emph{som Aarsag. Fra Religionens og Troens Synspunkt
er al Virkelighed, den naturlige ligesaa
vel som den overnaturlige, i sidste og høieste
Instans at henføre under Aandsvirkelighed},
\emph{Virkelighed for Aanden}.

Dersom Virkelighedens eneste ubedragelige Kjendemærke
var det udvortes Sandselige, det ret Haandgribelige,
saa vilde Jordarter og Stene vel omtrent være at
sætte paa øverste Plads i Virkelighedens Rige. Vi fik da
en Materialisme af den allergroveste Art; thi Materialisterne
pleie dog ellers at lade alt Virkeligt beroe paa en
Vrimmel af Smaadele (Atomer), der ikke ere ligefrem
haandgribelige. Fornuft og Religion stemme deri overens,
at der ikke blot virkes i det Legemlige, men ogsaa
i det Ulegemlige, og at der som Følge heraf ligesaa vel
maa gives en ulegemlig som en legemlig Virkelighed.
De, der indrømme, at der er Fornuft i Tingene, maae da
ogsaa indrømme, at alt Legemligt har sin Grund i et
Ulegemligt, med mindre man da vil ansee den evige
Fornuft selv for et Legeme, en Art usynlig Materie.

Striden mellem Religion og Fornuft kommer først
tilsyne, naar Virkelighedens usynlige Aarsag nærmere
skal bestemmes. Fysiken holder sig til Materiens Kræfter
og forsker efter Love; den har kun med Virkelighedens
Naturside, altsaa kun med Naturvirkeligheden, ikke med
Aandsvirkeligheden, at gjøre. Filosofien bestræber sig for
at fremstille den virkelige Natur, den i sig sammen%
% --- p. 9 ---
\setcounter{footnote}{0}%
\opage{9}hængende Naturvirkelighed, som begrundet i Aandens
Væsen; men saaledes, at Aandens Virken i Naturen ikke
paa noget Punkt overskrider Naturbetingelsernes Grændse.
Religionen --- vi tale her om den aabenbarede Religion
--- indskærper os derimod som allerførste Grundsætning,
at Guds frie, almægtige Villie er den oprindelige Aarsag
til al Virkelighed. \emph{I Begyndelsen skabte Gud
Himlen og Jorden}; saadan lyde, som bekjendt, de første
Ord i den hellige Skrift; hermed er da udtrykkelig sagt,
at hele Verdens Virkelighed i Grunden er Aandsvirkelighed.
At troe paa den første Artikel om Skabelsen og
dog nægte, at Aandsvirkelighed er det sande omfattende
Udtryk for al Virkelighed, er en Modsigelse.

Forholdet mellem Natur og Skabning kan næppe betegnes
tydeligere end ved at sige: Skabningen i sin
Grund og Sammenhæng er \emph{Natur}; Naturen i sin Oprindelse
fra og Afhængighed af Gud er Skabning\footnote[1]{Det gjælder her, som over alt, om at holde Synspunkterne skarpt ud fra hinanden. At forklare det Naturlige halvt om halvt overnaturlig og det Overnaturlige halvt om halvt naturlig, er vrøvlagtigt; ved at sammenblande Fysik og Religion og mediere Blandingen i dogmatisk Spekulation gjør man Religion og Videnskab en slet Tjeneste.}. Bekræftelsen
paa, at de skabte Ting have deres Natur, er
at see i Lovene; Tegnet paa, at den virkelige Natur oprindelig
er skabt, er at see i Underet. Underet er enfoldigt
at bestemme som Almagtens Tegn, en udtrykkelig
Tilkjendegivelse af, at Gud virkelig har en fri Villie, og
at denne Villie ubetinget har alle Ting i sin Magt.
Dersom den virkelige Natur ikke var modtagelig for Miraklet,
saa var Naturriget ikke skabt: at nægte Underets

% --- p. 10 ---
\setcounter{footnote}{0}%
\opage{10}Mulighed er at nægte den frie Skabelses Mulighed. Det
virkelige Under, hvad enten det forøvrigt foregaaer i det
Indvortes eller i det Udvortes, i det Ulegemlige eller i
det Legemlige, er da i allerstrengeste Forstand at kalde
Aandsvirkelighed, Virkelighed for Aanden.

Den hellige Skrift stiller overalt, hvad Virkelighed
angaaer, Miraklerne i det Aandelige paa lige Trin med
Miraklerne i det Legemlige; dette er et Hovedtræk, som
endelig ikke maa oversees. Nogle Exempler. Jehova
aabenbarer sig for Abraham, først i et Syn (1. Mos. 15.
Kap.), dernæst (18. Kap.) i en udvortes virkelig Skikkelse;
den første Aabenbarelse viser altsaa hen paa det Ulegemlige,
den sidste paa det Legemlige. Hvo vil nu fra
Troens Standpunkt paastaae, at den første Aabenbaringsmaade
er mindre overnaturlig, mindre virkelig og derfor
mindre tilforladelig end den sidste? At den sidste gaaer
stærkt ind paa det udvortes Naturlige, sees af den Gjæstfrihed,
hvormed Patriarken beværter de tre Himmelske:
»han tog tyk Mælk og sød Mælk og den Kalv, som han
havde tillavet, og satte den foran dem, og selv stod han
foran dem under Træet, og de aad«. Gaaer det nu an
at sige, at den første Virkelighed, Virkeligheden i Synet,
kun er en Smule Aandsvirkelighed, en flygtig Ide-Virkelighed,
maaske en vaagen Drøm, en Fantasi af Fader
Abraham, men at den sidste er den egenlige, rette og
solide Virkelighed, fordi det er en Virkelighed med Kalvekjød
og tyk Mælk, altsaa en Virkelighed, der falder udenfor
Aanden?

Der er i Henseende til Virkelighedsbegrebet selv
ingen Forskjel mellem det Gamle og det Nye Testamente.
Vi lære det af Opstandelsens Under. At Kristus virkelig
% --- p. 11 ---
\setcounter{footnote}{0}%
\opage{11}er opstanden fra de Døde, at hans Aabenbarelse for Sine
ikke blot er skeet \emph{indvortes} som i et Syn, men ogsaa
\emph{udvortes}, ja (Joh. 20, 27; jvnfr. Mtth. Mr. Luc.) paa
ret \emph{haandgribelig} Maade, er saa kraftig bevidnet, som
det vel er muligt, af troende Vidner; men Ingen vil dog
vel nægte, at det Legeme, hvormed Kristus blev korsfæstet, og hvormed han opstod, er ved Opstandelsen undergaaet
en saadan Forvandling, at han snart er kjendelig,
snart igjen ukjendelig, at han kommer som et Syn og
dog som den, der har Kjød og Ben og Sands for Næring
--- de gave ham et Stykke af en stegt Fisk og af
en Honningkage (Luc. 24, 42) --- at, med andre Ord,
hans legemlige Virkelighed har optaget det Naturlige i
det Overnaturlige, at den i Et og Alt svarer til Aandsvirkelighed,
Virkelighed for Aanden. De Theologiske, som
ikke kunne tilfredsstilles ved det Aandsvirkelige, men
nødvendig maae have enten »Honningkage« og »stegt
Fisk« eller »Kalvekjød« og »tyk Mælk« i Tilgift for at
bøde paa det Aandelige og udfylde Virkeligheden, skulle
med al deres Theologi have ondt nok ved at troe paa
Kristi virkelige Opstandelse.

\begin{center}III.\end{center}

\emph{Aandsvirkelighed er det kategoriske Udtryk}
for \emph{Enheden af Aabenbaringens Form med det
aabenbarede Indhold. Det er den paa een Gang
overnaturlige og naturlige Enhed af Aanden og
Ordet, der gjør Ordet til Guds Ord og Gudsordet
selv til høieste Virkelighed, de virkelige}
Kjendsgjerningers Bærer og Fremstiller. Det
\emph{fremstillende Ord har ikke den paagjældende
% --- p. 12 ---
\setcounter{footnote}{0}%
\opage{12}Virkelighed udenfor sig, men i sig.} Skilles
\emph{Ordet fra sit aandsvirkelige} Indhold, saa \emph{bliver
det Sproglige, det Grammatiske Hovedsagen:
Det aabenbarede Ord forvandles til en
filologisk Text, og der opstaaer Strid mellem}
Aand og \emph{Bogstav}.

»I Begyndelsen var Ordet, og Ordet var hos Gud,
og Ordet var Gud« (Joh. 1, 1). --- Var det Opgaven,
kort og bestemt at angive Synspunktet for den glosepillende,
stokfilologiske Exegese, der gjør Gudsordet til
Sætninger og Sætningslæren til Grammatik, da kunde
det allerbedst skee ved passende Omskrivning af det anførte
Sted: »I Begyndelsen var Grammatiken, og Grammatiken
var hos Gud, og Grammatiken var Gud« o. s. v.
Fra et objektivt lærd Standpunkt vil det næppe falde
Nogen ind at kaste Vrag paa et dybtgaaende filologisk
Studium af Grundtexten i Gl. og N. Test.; heller ikke
skal det oversees, at Udbyttet af et saadant Studium
middelbart afgiver en Veiledning ogsaa for dem, som
alene i religiøs Interesse ville sikkre sig Forstaaelsen af
den bibelske Text. Men der maa i Religionens Navn
nedlægges Protest mod den minutiøs-filistrøse Vigtighed
og stadig tiltagende Vidtløftighed, hvormed Bibelfilologien
i den nyere Tid har vidst at gjøre sine exegetiske Analyser
til en videnskabelig Troesfornødenhed, et Slags
grammatikalsk Salighedsanliggende. Saa skeer det, at
Bisagen bliver Hovedsag, og Hovedsagen Bisag; saa skeer
det, at Gudsordet løses op i Sætninger og Bogstaver.

Forholdet mellem Gudsordet og den bibelske Text
kan i Korthed oplyses paa følgende Maade. I Opstandelsesevangeliet
eier den Troende en fuldkommen sand og
% --- p. 13 ---
\setcounter{footnote}{0}%
\opage{13}tilforladelig Fremstilling af en virkelig historisk Begivenhed.
Opstandelsens Under er Guds Gjerning; den ufeilbare
Beretning om dette Under er selv et Under, er
ogsaa en Guds-Gjerning: Begivenheden og dens Opfattelse,
den udvortes og den indvortes Virkelighed smelter i Gudsordet
sammen til en eneste udelbar Virkelighed, Aandsvirkelighed.
Og dog er der Forskjel mellem Beskrivelsen
og det, der beskrives. Den evangeliske Beretning er et
Speil, der gjengiver Virkelighedens enkelte Træk i et tro
Billede. Troens Øie seer i Speilet paa Billedet og i
Billedet paa Virkeligheden, den virkelige Begivenhed:
Blikket er forsaavidt udadvendt. Dersom der nu i udvortes
Forstand kun var eet Evangelium, een Beretning,
saa vilde det udadvendte Blik uvilkaarlig blive indadvendt,
Enheden og Harmonien mellem Kjendsgjerningen og Opfattelsen,
mellem det udvortes og det indvortes Virkelige,
vilde da holde sig uforandret. Men den guddommelige
Sag har og maa i Kraft af Aabenbaringens Væsen nødvendig
have sin menneskelige Side. Det ene Evangelium
er delt i fire Evangelier, den ene ufeilbare Beretning i
fire indbyrdes forskjellige og dog ufeilbare Beretninger.
Forseer man sig paa denne Forskjel, saa fremkommer
der en udvortes Adskillelse mellem Ordet i guddommelig
Sandhed og Beretningen i menneskelig Form, en betænkelig
Skilsmisse mellem Gudsordet og den bibelske
Text; der bliver Noget at bestille for Grammatiken.

Hos Matthæus (Kap. 28) begynder det med, at Maria
Magdalene og den anden Maria kom, da det gryede ad
den første Dag i Ugen, for at see til Graven. Hvilket
Syn! Herrens Engel farer ned fra Himlen at vælte Stener
fra Graven; hans Skikkelse er som Lynet; Vagten flygter
% --- p. 14 ---
\setcounter{footnote}{0}%
\opage{14}og Engelen tiltaler Kvinderne. Men Ingen af de tre
andre Evangelister synes at vide Noget om, at der var
sat Vagt ved Graven; heller ikke har, ifølge deres Beretning,
Nogen iblandt Kvinderne været Øienvidne til, hvorledes
Stenen blev bortvæltet; tvertimod, det var allerede
sket før de kom. Endnu en Uoverensstemmelse. Markus
(K. 16, V. 1) siger udtrykkelig, at Kvinderne først, da
Sabbaten var forbi, gjorde Indkjøb til Balsameringen;
men Lukas (Kap. 23, V. 56) bevidner, at de endnu før
Sabbaten beredte Urter og Salver: hvorledes kunde de
tilberede de Urter og Salver, de endnu ikke havde indkjøbt?
Om man ogsaa ved filologisk Behandling af
Verbet: »\emph{beredte}« (ἡτοίμασαν) kan glide nogenledes
hen over denne tilsyneladende Modsigelse, hvorledes slipper
man da igjennem Vanskelighederne med de modsigende
Beretninger om Englene? Matthæus og Markus
bevidne, at der kun var een Engel ved Graven, men
Lukas og Johannes anføre udtrykkelig, at der var to, og dog
komme deres Vidnesbyrd heller ikke udvortes overens. Ifølge
Lukas blive Englene sete af alle »Kvinderne fra Galilæa«;
men ifølge Johannes faaer ingen Anden dem at see end
Maria Magdalene. Endnu mere paafaldende er den forskjellige
Maade, paa hvilken Englene udtale sig. Hos
Lukas holde de et udførligt, belærende Foredrag for
Kvinderne; hos Johannes indskrænker Tiltalen til Maria
Magdalene sig til det Spørgsmaal: Kvinde, hvi græder
du? Ifølge Matthæus og Markus byder Kristus umiddelbart
efter sin Opstandelse, saavel ved Englen i Graven
som ved sit eget mundlige Ord til Kvinderne, at Disciplene
skulle møde ham i Galilæa; men ifølge Lukas befaler
han, at de alle skulle forblive i Jerusalem, indtil
% --- p. 15 ---
\setcounter{footnote}{0}%
\opage{15}den Helligaand var udgydt over dem. Ved saaledes at
gaae i det Enkelte vil Læseren opdage den ene paafaldende
Uoverensstemmelse efter den anden. Lessing har
i Forsvaret for sin »Ubekjendte«\footnote[1]{\emph{Eine Duplik}. 1778. Sämmtliche Schriften, 4ter Band. Berlin 1825. S. 95 flgd.}, paavist ikke mindre
end »ti Modsigelser«. Ti Modsigelser i fire korte Beretninger!
er det ikke forfærdeligt?

Dersom det i slige Modsigelser Anstødelige virkelig
havde sin Grund i Skriften selv, saa burde en saadan Skrift
ikke bære Navnet: den hellige Skrift; det vilde være mere
passende at kalde den den dunkle, forvirrede, modsigende
Skrift. Men Sagen er, at den ved Modsigelserne vakte Forargelse
ikke har sin Grund i Skriften, men i Skriftfortolkernes
falske Udgangspunkt. Ved at begynde med en filologisk
Kritik af Texten faaer man Troen uforvarende narret ind paa
Videns Omraade og kommer til at anlægge en Maalestok,
der staaer i aabenbar Strid med Aabenbaringens Væsen.
Grundvildfarelsen er, at man absolut vil behandle Bibelen
som ethvert andet skriftligt Oldtidsdokument og ansee det
for en afgjort Sag, at der i den hellige Skrift lige saa
fuldt som i profane Skrifter kan og maa gjøres en endelig
kritisk Adskillelse mellem de i Ordet fremstillede
Kjendsgjerninger og disse Kjendsgjerningers ydre formodede
Virkelighed. Men i den hellige Skrift er det
netop en Selvmodsigelse at gjøre denne Adskillelse. Underet
og Troen ere uadskillelige; tag Underet bort, og der
er ingen Tro; tag Troen bort, og der er intet Under.
Det er ikke blot ved den senere Nedskrivning, at det
væsenlig kommer an paa Troesvidnets eiendommelige Op%
% --- p. 16 ---
\setcounter{footnote}{0}%
\opage{16}fattelse og Fremstillingsgave; nei Kjendsgjerningen og
dens Opfattelse ere fra første Øieblik af i Spiren selv
uadskillelig forenede: den udvortes Virkelighed i Begivenheden
og den indvortes i Opfattelsen er een udelelig
Kjendsgjerning, een Virkelighed: det er Aandsvirkelighed,
Virkelighed for Aanden.

At Sagen virkelig forholder sig saaledes, sees allertydeligst
deraf, at de Biomstændigheder, der fra et kritisk
filologisk Synspunkt skulde tale \emph{imod}, netop komme til
at tale for Evangeliets Troværdighed. a) Elementerne af
det Udvortes tilhøre alle een Virkelighed, een stor, vidunderlig
Begivenhed. Men Begivenheden har flere Led
og mange Enkeltheder; disse elementære Enkeltheder ere
saaledes opløste i Tilegnelsen og forvandlede i Fremstillingen,
at det udvortes og indvortes Oplevede ikke lader
sig adskille i Ordet. Ordets uopløselige Enhed af det
Udvortes og det Indvortes er overalt et Karaktermærke
for den aabenbarede Sandhed. b) De Afvigelser i Opfattelsen,
som, naar man stirrer mod det Udvortes og
hænger sig i Bogstaven, blive til virkelige, anstødelige
Uoverensstemmelser, forvandle sig til talende Vidnesbyrd
om levende, fri, personlig Tilegnelse, naar de sees i Aanden
og forstaaes i Ordet. Uoverensstemmelser i en Fortælling
med andre Fortællinger kunne unægtelig frembyde
haarde Knuder, uopløselige Modsigelser, naar den Fortællende
giver sig ud for at kopiere en udvortes given
Virkelighed; thi det er bogstavelig sandt, at een Engel
er ikke to; at det, man kjøber før Sabbaten, kjøber man
ikke efter Sabbaten; at Galilæa ikke er Jerusalem o. s. v.
Men det er i Aanden og Ordet endnu mere afgjørende
sandt, at i Forhold til en overvældende stor og mirakuløs
% --- p. 17 ---
\setcounter{footnote}{0}%
\opage{17}Begivenhed maa hvert enkelt Vidne være saa stærkt
grebet, saa ganske opfyldt og inderlig bevæget af Begivenheden
i dens Helhed, at Tilegnelsens menneskelige
Side ogsaa vil lade sig tilsyne, saa at han i Gjengivelsen
af Hovedpartierne og under en subjektiv Udmaling af de
enkelte Træk uvilkaarlig kommer ind paa individualiserende
Ændringer. Evangelisten er bevæget af Aanden
og i Aanden af Virkeligheden; men han er ingen Fotograf,
der kun vil give os en Kopi af en udvortes Gjenstand,
saa lidt som han er en Skriverkarl, der mekanisk nedskriver,
hvad Aanden mekanisk dikterer. Evangelisten er
et Sandhedsvidne, forstaaer sig, ikke Sandhedsvidne i dogmatisk-bispelig-moderne
Forstand --- til en saadan Aandrighedens
Høide kan han naturligvis ikke hæve sig ---
men han er dog virkelig et Sandhedsvidne; han vidner
i Aanden om Virkelighed: Vidnesbyrdets Indhold i Ordet
er Aandsvirkelighed.

\begin{center}IV.\end{center}

\emph{Aandsvirkelighed er betegnende for et
Standpunkt, hvorfra det bliver muligt} at \emph{gjennemføre
Skilsmissen mellem den videnskabelige
og den religiøse Opfattelse af den hellige
Historie saa principielt, at der saavel om det
Hele som om dets enkelte Dele kan danne sig
to Slags Kritik: Videnskritik og Troeskritik}.

Blandt mange dels overfladiske, dels underfundige
Vendinger, hvorunder de paa Viden Troende have søgt at
afvise den strenge Modsætning mellem Religion og Videnskab,
hører ogsaa den, at en religiøs Fortolkning af den
hellige Skrift maa blive grundløs, vilkaarlig og almueagtig,
% --- p. 18 ---
\setcounter{footnote}{0}%
\opage{18}naar den løsrives fra videnskabelig Forskning. Det er,
siger man, at kaste Skygge paa Kundskab og Lærdom,
naar man bilder det menige Folk ind, at enfoldigt
Troende kunne forstaae Evangeliet med det bare Hjerte
uden Veiledning af virkelig Skriftklogskab. At lade Alt
komme an paa en levende Tro og gaae ud fra den simple
Paastand, at Mirakler ere Mirakler, at man enten maa
antage dem alle eller forkaste dem alle, er en altfor bekvem
Maade at slippe fra Vanskelighederne paa; kort;
Theologien vil imponere Almenheden ved at paakalde
Grundigheden. Vi maae hertil bemærke, at hoverende
Basunstød i Grundighedens Trompet kun ere Nødsignaler,
der Intet have med Aandsvirkelighed at gjøre. Vi ere
saa lidt tilsinds at ville enten nedsætte Grundigheden
eller smigre den selvkloge Tankeløshed, at vi tvertimod
efter ringe Evne ere betænkte paa at komme de Oppustede
tillivs og forfølge Principløsheden til det Yderste, enten
den forøvrigt stammer fra Skriftklogskab eller fra Selvklogskab
eller maaske fra begge i Forening.

Det bliver med Rette fastholdt, at der maa gives en
videnskabelig Fortolkning af den hellige Skrift ved Siden
af den religiøse; men det skal atter og atter gjentages,
at de to Fortolkninger tilhøre hver sin Erkjendelseskreds,
at den ene ikke kan begrunde den anden, men at de vel
kunne udøve en middelbar, indirekt Indflydelse paa hinanden\footnote[1]{Det forudsættes naturligvis, at de Læsere, som ønske at sætte sig ind i den her udtalte Synsmaade, efter Tid og Leilighed gjøre sig bekjendte med de herhen hørende Skrifter: \emph{Hindringer og Betingelser for det aandelige Liv}, \emph{Den gode Villie som Magt i Videnskaben}, \emph{Religionsfilosofi}, Afhandlingen om \emph{Aandsvirkelighed} (Tidsskriftets 10de Hefte, Oktober 1870).}. Det er nu dette Forhold, vi skulle oplyse ved
% --- p. 19 ---
\setcounter{footnote}{0}%
\opage{19}en kort Karakteristik af den til Viden og den til Troen
svarende Kritik.

a) \emph{Videnskritiken} karakteriserer sig ved Spaltning
og Udstykning af det i Aandsvirkelighed sammensluttede
Hele. Kjendsgjerningen falder udenfor Opfattelsen; Beretningen
er Et og den virkelige Begivenhed et Andet:
en Overensstemmelse er mulig, men --- denne Tvivl er
Kritikens Spore! --- en Uoverensstemmelse er jo ogsaa
mulig. Ved omhyggelig Sammenligning af enhver foreliggende
Beretnings enkelte Træk med hverandre indbyrdes
og med andre tilsvarende Beretninger, hvis saadanne
gives, stræber Fortolkeren at komme den bagved liggende,
ubekjendte Virkelighed saa nær som muligt. Kritikens
Særkjende er, at Forskeren overalt gaaer uden om Aanden
og bagved Ordet, for at lede efter Kjendsgjerningen,
det udvortes Virkelige. Hvad er da virkelig skeet? Er
Kristus virkelig blevet fristet af Djævelen? Er han virkelig
blevet forklaret paa Bjerget? Har han virkelig vandret
paa Havet? Mon han virkelig ved sit Ord har formaaet
at stille Stormen? o. s. v. Nu gjælder det om Grundighed;
jo grundigere den filologisk-kritiske Fortolker gaaer
til Værks, desto tydeligere viser sig Kritikens Udbytte.

Udbyttet er, at Virkeligheden forsvinder i et mystisk
Halvmørke; de naturlige Træk blive staaende, men de
overnaturlige skimtes kun i »Sagn og Myther«\footnote[1]{Jvnfr. En Fortolkning af den hellige Historie (Tidsskr. 4de Hefte, April 1870).}.

\emph{Fristelsen i Ørken} har ingen Virkelighed, hverken
for Djævlens eller for Kristi Vedkommende. »For

% --- p. 20 ---
\setcounter{footnote}{0}%
\opage{20}Djævlens Vedkommende maa Fremgangsmaaden ved den
hele vanskelige Kamp, det Platte i Ordnen og Indleden
af Anfægtelserne synes ligesaa vanskeligt at forene med
den Kløgt og beregnende Snedighed, der her dog skulde
naae sit Høidepunkt, som, fra Jesu Side, Taaligheden i
at fortsætte disse Forhandlinger, og af den onde Aand at
lade sig føre fra Sted til Sted gjennem Prøvelsens forskjellige
Stadier, vanskeligt forenes med Idealet af den
sædelige Renhed, der strax maatte have draget sig ud af
Berøring med det Urene«. (H. N. Clausen.)

»\emph{Forklarelsen paa Bjerget}« \emph{har ingen sand}
Virkelighed; her er --- trods den stærke Belysning ---
kun »et Taagebillede i mystisk Halvlys«. »De trende
Øienvidner (Petrus, Jakobus og Johannes) vare betyngede
af Søvn, de vidste ikke, hvad de talede, thi de vare forfærdede;
da de ret kom til Besindelse og saae sig om,
saae de Ingen uden Jesus alene». (H. N. Clausen.)

»\emph{Vandringen paa Søen}« viser jo vistnok hen til
en virkelig Sø og en virkelig Vandring, men det i Evangelierne
betegnede Forhold mellem Sø og Vandring er
neppe virkeligt. Den filologiske Kritik har »grundet
Formodning om, at Episoden om Peder er en senere Tilsætning,
der skyldes den græske Oversættelse eller Bearbeidelse
af Matthæus-Evangeliet, en Formodning, som
bestyrkes ved den hele Karakter af denne Fortælling, der
minder om den udsmykkende og forstærkende Tendens i
de apokryfiske Evangelier«. (H. N. Clausen.)

»\emph{Stormens Stillen paa den galilæiske Sø}«
har i den evangeliske Fremstilling i den Grad faaet sin
Virkelighed forvreden, at Kjendsgjerningen selv er bleven
vendt baglænds. Evangelisterne (Mtth., Mrk., Luk.) have
% --- p. 21 ---
\setcounter{footnote}{0}%
\opage{21}nemlig i deres Enfoldighed antaget, at det var Stormen,
som Jesus truede, medens Kritiken maa finde, at det
ikke er Stormen, men Disciplene, der faae en Revselse;
thi det er langt rimeligere at antage, at Jesu truende
Ord »har været henvendt til Disciplene end til Elementerne,
til Veir og Vind«. (H. N. Clausen.)

Det er ikke blot et vist indskrænket Tal af enkelte,
for den religiøse Bevidsthed maaske mindre væsenlige
Undergjerninger, hvis Virkelighed paa denne Maade
bliver opløst og omkalfatret, thi de store Frelsens Kjendsgjerninger,
de Grundundere, paa hvilke den hele Aabenbaring
hviler, friste samme Skæbne.

At en Kritik, der gaaer udenom Ordet og bagved
Aanden, ikke kan anerkjende Kristi Undfangelse af Aanden,
er en Selvfølge. Hvad skulde vel bevæge den til
en saadan Anerkjendelse? Forevangeliets »fragmentariske
Fortællinger« ere af tvivlsom Oprindelse; »en mythiserende
Indklædning«, en »udsmykkende Indfatning«, er
traadt »udfyldende og formende« til for at bøde paa »det
Sparsomme og Vaklende« i de historiske Overleveringer.
--- Med Slutningen gaaer det akkurat som med Begyndelsen.
Evangelierne om Kristi Opstandelse og Himmelfart
ere ligesaa mythiserende, ligesaa fragmentariske,
ligesaa udsmykkende som Evangelierne om hans Fødsel
og Barndom. »De indbyrdes Afvigelser, med hvilke de
fire evangeliske Beretninger om Opstandelsen gjennemkrydse hverandre«, virke i den Grad forstyrrende, at »der
kun med stor Usikkerhed lader sig, ved den indbyrdes
Sammenstilling, tilveiebringe en nogenlunde sammenhængende
Fortælling«. Endnu bedrøveligere gaaer det
med Himmelfarten. At den bevidnes som et Grund%
% --- p. 22 ---
\setcounter{footnote}{0}%
\opage{22}faktum, maa indrømmes; men hvori dette Grundfaktum
egenlig har bestaaet, er trods Granskningens Grundighed
ikke til at udgranske; at sige med rene Ord: Evangelisterne
vide ingen Besked, og Fortolkerne ikke heller.

b) \emph{Troeskritiken} er nu kjendelig paa, at den holder
Kjendsgjerningerne fast i den evangeliske Opfattelse.
Det er uimodsigelig vist, hvad Pastor Birkedal fremhæver,
at det troende Hjerte ikke kan nøies med »Skyggegjerninger»,
men trænger til »virkelige, haandgribelige Kjendsgjerninger«;
men hvorledes den høitærede Forf. er falden
paa den Tanke, at jeg muligvis kun anseer Aandsvirkelighed
for en Art Ide- Virkelighed og de aabenbarede Kjendsgjerninger
for Skyggegjerninger, maa han selv forstaae,
thi jeg forstaaer det ikke. Hvad jeg, ifølge Troeskritikens
Princip, ufravigelig maa fastholde, er, at den Troende med
al sin Trang til »haandgribelige Kjendsgjerninger« ikke
tør gaae bagved Ordet og udenom Aanden; han maa
ikke bilde sig ind, at han kan see Evangelisten over
Skulderen og ved at kige som ud af et Vindue efter et
Udvortes faae Øie paa en bedre, sikkrere og tilforladeligere
Opfattelse af Virkeligheden end den, der er ham given i
det evangeliske Ord.

Lad os nu see, hvorledes denne for Aandsvirkeligheden
gjældende Hovedsætning kan anvendes. Pastor B.
giver Sagen følgende Vending. »Hjertet siger: Herrens
Jærtegn maae være virkelige. --- Han maa virkelig have
kunnet byde over Verdenskræfterne og ikke behøve at
lade sig binde af Naturlovene --- det er nødvendigt for
min Tro; thi kunde han ikke gjøre disse Undere i Naturen,
da kan jeg ikke troe paa ham selv som det største
Jærtegn, som den, der er Naturens Herre, da kunde jeg
% --- p. 23 ---
\setcounter{footnote}{0}%
\opage{23}heller ikke troe paa, at Han gjør det endnu større Under,
trods min Natur og dens Love, der efter \emph{deres}
Nødvendighed maatte føre til Døden, at give mig sit
Liv og lade mig med sig gaae over Dødshavet, saa jeg
ikke gaaer tilbunds deri. Det har jeg paa \emph{Følelsen}:
kunde han ikke bringe Naturen der udenfor til at lystre,
da kunde han heller ikke bringe min endnu mere oprørske
Natur til at underlægge sig Ham«. Dette er i
Troens Navn baade vist og sandt i alle Maader. Men
dersom nu Pastor B. mener, at han »paa Følelsen« kan
have en sikker Viden om, at Kristus virkelig har vandret
paa Havet, at han ved sit Almagtsord virkelig har stillet
Stormen o. s. v., da beder jeg ham betænke, om han
ikke her har tilladt sig en forhastet Slutning. Saalænge
han i al Almindelighed taler om, hvad Frelseren maa
have \emph{kunnet} for at kunne være Frelser, giver jeg ham
ubetinget Ret; thi han taler om en Mulighed; det er
en sand Troeslogik, hvoraf han giver et Udkast. Men,
naar han gaaer i det Enkelte og vil »paa Følelsen« afgjøre,
hvilke Undergjerninger, Frelseren \emph{virkelig} har gjort,
kan jeg ikke give ham Ret, thi fra Mulighed til Virkelighed,
fra det Logiske til det Antilogiske, er der et
Spring.

Har Kristus virkelig vandret paa Havet? Det er der
intet Menneske, troende saalidt som ikke-troende, der
veed det, thi om Sligt gives der ingen Viden. Det er jo
muligt, at Kristus ikke har gjort det, skjøndt han vel
har kunnet gjøre det. Sæt, at han ikke har vandret paa
Havet, men derimod virkelig stillet Stormen, er han saa
ikke Frelseren? Det er vist, at den Samme, der har
% --- p. 24 ---
\setcounter{footnote}{0}%
\opage{24}kunnet opvække Lazarus i Bethanien\footnote[1]{Pastor Møller lægger særlig Vægt paa Miraklet med Lazari Opvækkelse, og forudsætter med Grund, at dette Mirakel efter min Synsmaade maa have „fuld Virkelighed“. Ja, det maa have fuld Virkelighed! Dets Virkelighed er for alle Troende bekræftet ved saa stærke Vidnesbyrd, som det vel er tænkeligt. Men at det i denne Henseende skulde have noget særeget Fortrin fremfor saa mange andre Jesu Mirakler, kan jeg ikke indrømme. Ordene: „stinker allerede“, kunde jo af Tvivlere let opfattes som en blot Formodning af Martha, støttet paa, at Lazarus allerede havde „ligget fire Dage“ i Graven. Dertil kommer, at Nogle blandt de Tilstedeværende jo gik hen og gav de Ypperstepræster en Underretning, som var af den Beskaffenhed, at Raadet derefter besluttede, at Jesus skulde døe. Men at ville true en Mand med Døden, som vitterlig for Alles, ogsaa for de vantroe Tilskueres, Øine har formaaet at kalde et allerede stinkende Lig tillive, er ikke blot Tegn paa Vantro, men umiskjendeligt Tegn paa en til Daarekisten kvalificeret Afsindighed. Man betænke dog, hvilket Indtryk, en saadan Gjerning maatte gjøre, naar den, vel at mærke, var paa udvortes fysisk Maade bleven en hel Forsamling vitterlig; alene Efterretningen derom maatte jo slaae det mest forstokkede Synedrium med Rædsel. Efter min Synsmaade --- det er et Punkt, jeg ret vil anbefale til troende Mænds alvorlige Eftertanke --- kan intet Mirakel i Jesu Liv være skeet paa en saa fysisk iøinefaldende Maade og under saa udvortes overskuelige Omstændigheder, at ingen sadducæisk Tvivler skulde fundet Udvei til en naturlig Forklaring. At gjøre et Mirakel saa aabenlyst og paa en saadan Maade, at der kun behøves Øine og Øren, Sands og Forstand, men ingen Tro, for at forvisse sig om dets Gyldighed som Mirakel, er at ophæve Miraklets religiøse Betydning. Det var netop en saadan Misbrug af det Hellige, Djævlen vilde fremkalde, da han foreslog Kristus at styrte sig ned fra Templets Tinde --- et Experiment i det udvortes Virkelige, vistnok til Overraskelse og Forbauselse for et jerusalemsk Publikum.\par Den negative Kritik har ofte gjentaget den Indvending, at saa mange og saa store Undergjerninger som de, om hvilke Evangelisterne fortælle, umulig kunne være virkelig skete, uden at Herodes og Pilatus og Keiseren i Rom maatte have hørt derom og ladet Sagen juridisk undersøge. Denne Indvending falder til Jorden, naar vi erindre, hvad det egenlig vil sige, at Miraklets Virkelighed kun er til for Troende. --- Grundene for denne min Synsmaade har jeg med særligt Hensyn paa Lazari Opvækkelse udførligere fremsat i Efterskriften til »Forelæsningerne over Hindringer og Betingelser for det aandelige Liv i Nutiden«.}, maa ogsaa have
kunnet opvække Enkens Søn i Nain; sæt, at Kristus
% --- p. 25 ---
\setcounter{footnote}{0}%
\opage{25}ikke har opvakt Enkens Søn, men derimod Lazarus eller
omvendt: Enkens Søn, men ikke Lazarus, er han saa ikke
Frelseren? Det er en mislig Sag at ville i Tillid til Følelsen
indlade sig med den ræsonnerende Forstand om
Virkeligheden og den virkelige Beskaffenhed af de enkelte
Mirakler i deres Enkelthed.

Det, som vistnok, naar vi see nærmere til, er Pastor
Birkedals Mening, er dette, at der i Tillid til Evangelisternes
Troværdighed ud fra Følelsesindtrykket af det
Hele kan dømmes sikkert om det Enkelte. I Vandringen
paa Søen, i Herredømmet over Stormen, i Opvækkelsen
af Lazarus, i alle enkelte Mirakler, Optrin og Livsbilleder
gjenkjende vi overalt i alle Evangelierne den samme
Kristusskikkelse, den samme hjælperige Haand, den
samme Kjærlighedens Almagt: »det er Frelseren selv, det
ligner ham, Sligt har han gjort, saaledes seer han ud«.
Den umiddelbare Følelse af Kjendsgjerningernes Virkelighed
falder i alle Punkter sammen med en tilsvarende
Forvisning om, at Evangelisternes Opfattelse af denne
Virkelighed er ubedragelig sand; men naar Virkeligheden
ikke kan falde udenfor den evangeliske Opfattelse, altsaa
ikke udenfor Aanden, saa er det jo dog uimodsigeligt,
at den fra Først til Sidst maa bestemmes som Aandsvirkelighed.

Troeskritiken kan indrømme Ændringer, kronologiske
Omsætninger, Udvidelser, Forkortelser, ja ydre Uoverensstemmelser
i Fremstillingen af det Enkelte; men i disse

% --- p. 26 ---
\setcounter{footnote}{0}%
\opage{26}Omskrivninger er der ikke Spor af Forvanskning, ikke i
fjerneste Maade Antydning af nogensomhelst Misforstaaelse,
der kan vække Tvivl om Opfattelsens Gyldighed.
Sandhedskriteriet er, at hvert enkelt Livsbillede øiensynlig
svarer til det hele Grundbillede: er det Enkelte falsk,
saa er det Hele falsk; er det Hele sandt, saa er det Enkelte
sandt, og dermed Punktum.

c) \emph{Vexelforhold af Videnskritik og Troeskritik}.
Videnskritikens Fremgangsmaade er korrekt,
forsaavidt den lægger an paa en Tilnærmelse til objektiv
Vished i objektiv Viden; men den gjør en latterlig Figur
ved at tee sig, som om dens Tvivl og Uvished skulde
kunne gjøre den Troende uvis eller bevirke nogensomhelst
væsenlig Forandring i Troen. Lad det kun fra den objektive
Videns Synspunkt være uvist, hvorledes det virkelig
er gaaet til med Vandringen paa Havet; lad der i
Grammatikens Navn være »grundet Formodning om», at
»Episoden om Peder er en senere Tilsætning, der skyldes
den græske Oversættelse eller Bearbeidelse af Matthæus-Evangeliet«;
har denne »grundede Formodning« saa
mindste Indflydelse paa Troen? Ja, hvis der forelaa et
uigjendriveligt Bevis for, at Begivenheden ikke er virkelig
eller dog ikke saaledes virkelig, som den er fremstillet, det
var en anden Sag; saa vilde den Troende, efter hvad ovenfor
er sagt om Mulighed og Virkelighed, her i dette enkelte
Tilfælde ændre sin Forestilling om et udvortes Virkeligt.
Men kunde en saadan Ændring da i mindste Maade
rokke hans Tillid til Evangelistens Opfattelse af den
virkelige Kristus? Minder hvert Træk i Billedet ikke om
Frelseren som han er? Har ikke »Episoden om Peder«
noget i religiøs Forstand saa Gribende, at hver troende
% --- p. 27 ---
\setcounter{footnote}{0}%
\opage{27}Sjæl maa føle sig bevæget deraf? Sandelig, hvis Kristus
ikke har gjort, hvad der fortælles om ham i dette Evangelium,
saa \emph{kunde} han have gjort det, thi --- for at
tale Sammenligningens Sprog --- hans øvrige Gjerninger i
de øvrige Evangelier svare dertil. Altsaa, i værste Tilfælde,
vilde Kristusbilledet i det Hele forblive uforandret.
Men har Videnskritiken nu Vished om, at Begivenheden
ikke kan være virkelig? Ingenlunde. Men den har jo
dog »en grundet Formodning«. Ja, en Formodning, den
har faaet »grundet« ved at gaae uden om Aanden og
bagved Ordet. Skulde en saadan Formodning kunne
rokke Troen?

Forhandlingen mellem de to Slags Kritik vil uophørlig
dreie sig om Forholdet mellem Aand og Virkelighed.
Videnskritiken stræber fra Beretningen udad mod
de overnaturlige Kjendsgjerninger og naaer det Modsatte
af hvad den søger. Den vil gribe det Overnaturlige udenfor
Aanden og ender med en stykkevis Opløsning af alt
Overnaturligt. Den evangeliske Hovedbygning synker i
Grus, og de Virkelighedsstumper, der blive tilbage i al
Naturlighed, ere uden Betydning for Troen. Saa maa
dette da være en Spore for Troeskritiken til at gaae den
omvendte Vei. Ud fra andre Forudsætninger, med en
anden Psykologi, et andet Blik for Virkelighed og en anden
Maalestok for Mulighed, vil Troeskritiken stille sig to
Opgaver: 1) at tilegne sig de videnskabelige Oplysninger,
hvor de i Sandhed ere paa deres Plads og kunne gjælde
for virkelige Oplysninger, og derhos 2) at mærke
sig og blotte de Synspunkter, hvorfra Forstandsvildfarelserne
udgaae, Synspunkter, som, hvor fint end Traadene
ere spundne i det objektivt kritiske Væv, dog alle ere
% --- p. 28 ---
\setcounter{footnote}{0}%
\opage{28}kjendelige paa en af Viden halvveret Tro og en i Troen
forvrænget Viden. Til Fornuften i Videnskaben svarer
Naturvirkelighed; men til Aabenbaringens Aand svarer
Aandsvirkelighed.

\begin{center}V.\end{center}

\emph{Aandsvirkelighed er en Sammenslutning af
Extremer, hvis Midte afgiver Mulighed, ikke
blot for et Forlig mellem Tro og Viden}, men
\emph{ogsaa for en Adskillelse af Tro og Mening.}

Aabenbaringsvirkelighed, enten vi see paa Kjendsgjerningen
eller paa Opfattelsen, er Aandsvirkelighed; men
hermed er Synspunktet for Naturvirkeligheden ingenlunde
udelukket. Dette Synspunkt vil da ogsaa i den hellige
Historie gjøre sin Ret gjældende. Om Skabelsen, Paradiset,
Syndefaldet og Verdens Undergang er Videnskritiken
berettiget til at hævde Naturvirkeligheden en Betydning,
som den ikke har erholdt i den bibelske Fremstilling.
Idet Virkelighed staaer imod Virkelighed (Naturvirkeligheden
mod den i Aabenbaringsordet fremstillede),
beholder Viden sin Ret lige over for Troen; men den
herved fremkomne Strid udlignes i den Erkjendelse, at
Dualismen er sat af Aanden selv. Den samme almægtige
Guds Aand, der har skabt Verden, har ogsaa meddelt
Ordet om Verdens Skabelse. At Guds ubetinget frie
Villie er denne Verdens Aarsag, er den Sandhed, som
Aanden ved sit Aabenbaringsord vil indprænte Menneskeslægten
paa en for Alle fattelig, anskuelig Maade; Verdenskræfternes
uendelige Vexelvirkning, Lovenes Væsen
og Naturtingenes Sammenhæng skulde derimod ikke være
Gjenstand for Aabenbaring; de herhen hørende Under%
% --- p. 29 ---
\setcounter{footnote}{0}%
\opage{29}søgelser ere Fornuften og Videnskaben overladte. Videnskritiken
har altsaa Ret, naar den gjør opmærksom paa,
at Skabelsen i sex Dage ikke stemmer med den Naturudviklingens
Gang, som i Virkeligheden lader sig paavise.
Men den har Uret, naar den søger Grunden til Uoverensstemmelsen
deri, at den Aand, som taler igjennem
Skabelseshistorien, ikke har været inde i den nyere Fysik
og desaarsag viist sig uvidende om Tingenes virkelige
Sammenhæng. Troeskritiken har Ret, naar den gaaer
ud fra den Forudsætning, at den Aand, der udtaler sig i
Ordet om Skabelsen, er den almægtige, alvidende Aand,
den Faderens Kraft, der har frembragt Himlen og Jorden;
men den har Uret, naar den ved at laane lidt fra
Naturvidenskaben og konstle med det aabenbarede Ord
vil have Skabelseshistorien tillæmpet efter Astronomien\footnote[1]{Jevnf, Religionsfilosofien S. 151---66.}.   % PRINTED AS IS (Bind 5, p. 29): „Jevnf, Religionsfilosofien“ comma for the stop of Jevnf. (K M50, zoom z3)
Det er den theologiske Forplumring, den indbildte Enhed
af Tro og Viden, der danner Udgangspunktet for begge
Vildfarelser.

Paa Aabenbaringens Trin er den Aandens Virksomhed,
hvorved Ordet i Skabelseshistorien blev meddelt, en
ligesaa guddommelig Akt som den, hvorved Verden blev
skabt. Extremerne betyde, at Skaberakten i sig selv er
Mennesket ubegribelig, at derfor Aanden ved Ordet maa
oplyse de Troende, maa komme deres Fatteevne og religiøse
Trang imøde. Den Gud, der har aabenbaret sin
Almagt i Naturen, har ogsaa aabenbaret sin Villie i
Ordet: Aabenbaringen selv er Extremernes Midte.

Som det forholder sig med Skabelseshistorien, saaledes
forholder det sig med ethvert Afsnit i den hellige

% --- p. 30 ---
\setcounter{footnote}{0}%
\opage{30}Historie, hvor den i den fysiske Verdensorden udtalte
Virkelighed staaer i Modsætning til det aabenbarede Ord:
her er en væsenlig Modsætning, men ingen virkelig Strid.
Ved Opfattelsen og Bedømmelsen af den til Extremerne
svarende Midte aabner sig nu en Mulighed for Adskillelsen
mellem \emph{Tro og Mening}, en Adskillelse, der hidtil
desværre kun alt for meget, for ikke at sige ganske, er
bleven overseet i Religionslæren.

\emph{Troen paa Ordet i} Skabelseshistorien er og maa
være den samme, uforanderlig den samme, enten den
Troende forøvrigt hører til de Vidende eller de Ikke-Vidende.
Men den individuelle Opfattelse af Fremstillingens
Enkeltheder er ikke den samme. Den Ikke-Vidende tager
umiddelbart Alt efter Bogstaven uden Anelse om, at her
enten skal eller kan tænkes paa nogen Adskillelse mellem
Naturvirkelighed i de ydre Kjendsgjerninger og Aandsvirkelighed
i Ordet. Den Vidende derimod har gjort sig
Modsætningen klar; men han har ogsaa gjort sig det
klart, at der i den bibelske Fremstilling, enfoldig forstaaet,
ikke er Tøddel eller Bogstav, der i nogen Maade
kan svække Troen paa Ordets Guddommelighed. Saa
have de da begge, den Vidende og den Ikke-Vidende,
samme Tro, men det er langt fra, at de have samme
Mening.

Den kategoriske Forskjel mellem Tro og Mening er
let at paavise. Vi vælge, for at tage endnu et Par Exempler
med, Paradiset og Syndefaldet.

Der er nu ført saa megen ørkesløs Strid om den
geografiske Beliggenhed af Paradis, Beskaffenheden af
Kundskabstræet og den diaboliske Slange, at det vel
snart kunde være nok. De endeløse Kjævlerier falde
% --- p. 31 ---
\setcounter{footnote}{0}%
\opage{31}bort, naar man besinder sig paa Forskjellen mellem Tro
og Mening. At være enige i Troen det er: at være
enige i den religiøse Tilegnelse af Ordet, det er: at være
fuldstændig enige i Opfattelsen af det medskabte Gudsbillede,
Uskyldighedsidealet, Fristelsen, Faldet, Skylden og
Straffen, en Opfattelse, der kan og maa nøiagtig bestemmes
ved Fortællingens enkelte Træk og Ledføining. Det
er den religiøse Psykologis første afgjørende Forudsætninger,
hvorom det gjælder. Er der Enighed i dette
Stykke, saa er der Enighed i Troen. Hvad ydre Virkelighed
der kan tillægges Træerne, og hvorledes Slangen
har seet ud, derom kunne Meningerne være delte. De
individuelle Forestillinger, der ledsage Indtrykket af det
Hele, ere i den Grad betingede af de Enkeltes Følelse,
Fantasi og Dannelsestrin, at den Ene ikke tør paanøde
den anden sine Meninger.

Hvor vigtigt det er overalt paa de religiøse Anskuelsers
Omraade at gjøre grundig Forskjel mellem Tro og
Mening, vil da ogsaa være indlysende af den mangeaarige
Polemik for og imod det Grundtvigske Syn paa Troesbekjendelsens
Oprindelse~ og dens Forhold til Daabsordet,
der gaaer som den røde Traad igjennem vor nyeste
Kirkehistorie. »Livsnerven i den Grundtvigske Anskuelse«
maa, fra mit Udgangspunkt, udtrykkelig søges deri, »at
Troen paa Troesbekjendelsen i Principet falder sammen
med Troen paa dens Oprindelse fra Troens Gjenstand,
altsaa med Troen paa den i Bekjendelsens Ord personlig
nærværende Kristus. (Religionsfil. S. 72). Hvilke særlige
Forestillinger »Grundtvigianere« og »Anti-Grundtvigianere«
end ville sætte i Omløb angaaende det udvortes
Virkelige i den Art og Maade, paa hvilken Kristus i
% --- p. 32 ---
\setcounter{footnote}{0}%
\opage{32}sit »Kjøds Dage« skal have meddelt eller ikke meddelt
Apostlene Troesbekjendelsens Ordlyd, saa maa jeg dog
gjentage mit Enfoldige: det er der intet Menneske paa
Jorden, der veed; thi om Sligt gives der ingen objektiv
Viden. Saa bliver det da Opgaven at undersøge, hvad
der i denne Sag har afgjørende Betydning for Troen, og
hvad der kan betragtes som Meninger og være enten
rimelige eller urimelige Gjætninger. At Grundtvig ikke
har overseet Forskjellen mellem Tro og Mening, sees
blandt Andet af følgende Udtalelse: »Vi, som kalde os
Kristi Sendebud, og som jo skal meddele Menigheden al
den Oplysning i Troen, som Aanden ved Skriften eller
paa anden Maade skænker os, maae vel vogte os for at
gjøre vore Oplysninger enten om Troesbekjendelsen i det
Hele eller om dens enkelte Led og Dele til Troesartikler«\footnote[1]{Den kristelige Børnelærdom 1868. S. 45.}.

Hvilke Bestemmelser lade sig nu fastsætte som hørende
til Troen? Det kan jo kun være saadanne, som
uimodsigelig høre til Troens Grund og afgive de første
ufravigelige Betingelser for Troeslivet.

Daabsordet er et Sakramentord. Troen paa Sakramentordet
indeslutter den Forudsætning, at der ikke blot
døbes efter Kristi Befaling, men at Daabsakten selv i
Aanden ved Ordet er en Kristushandling\footnote[2]{Jevnf. Religionsfilosofien S. 508---514.}. Ordene:
Jeg døber dig i Navnet Faderens og Sønnens og den
Helligaands, maa da opfattes som Kristi Ord, Ord »af
Herrens egen Mund« til den, som lader sig døbe. Men
den, der siger: »Jeg døber dig«\dots{}. maa jo være den
% --- p. 33 ---
\setcounter{footnote}{0}%
\opage{33}samme som den, der forud har stillet Daabens Vilkaar
og modtaget Løftet om disse Vilkaars Opfyldelse gjennem
en oprigtig og frivillig Besvarelse af Spørgsmaalene:
»Forsager du\dots{}.? Troer du\dots{}.? Troer du\dots{}.?
Troer du\dots{}.? At lægge Kristus slige Ord i Munden,
og dog hemmelig nære den Biforestilling, at de egenlig
ikke ere hans egne Ord, men Ord af anden Oprindelse,
Ord, som med en vis velmenende Vilkaarlighed tillægges
ham, for at give Handlingen større Høitidelighed, hvad
er det? Det er at nedsætte Sakramentet til en Ceremoni,
det er at udhule Grundordet og begynde Troeslivet med
en Selvmodsigelse i Troen.

Opfattelsen af Sakramentet som en Kristushandling
i Menigheden fremgaaer med uafviselig Følgestrenghed
af Troens Princip, og Pastor C. Møller træffer nøiagtig
min Synsmaade, idet han minder om, at Intet lader sig
fastsætte om Troen uden det, som kan udledes af Principet.

Men nu har Troesbekjendelsens guddommelige Grundord
jo ogsaa en menneskelig Side. Om Sætningerne i de
tre apostoliske Artikler kan der, hvad Overleveringsmaade,
Oversættelsesudtryk o. desl. angaaer, naturligvis anstilles
filologiske og historisk-kritiske Undersøgelser, og der kan
igjen herom danne sig forskjellige Meninger. Hvorvidt
disse Meninger virkelig ere Troen underordnede, vil da
afhænge af, hvorvidt der fra alle Sider sigtes paa det
Oprindelige. I den aarvaagne og vedholdende Anstrengelse
for at sikkre og bevare Troesbekjendelsens Form i
alle enkelte Led og Udtryk svarende til det oprindelige
Troesord, det skabende Livets Ord, »Ordet af Herrens
egen Mund«, skulle vi da ikke see en ensidig Ivren
for en Partisag, men en oprigtig Kjærlighed til Sandhed.

% --- p. 34 ---
\setcounter{footnote}{0}%
\opage{34}At gjøre Troesbekjendelsen til et Slags ufuldstændigt
Kompendium af en »kristelig Dogmatik« og tage Indholdet
saa almindeligt, at man omtrent efter Smag og
Behag kan forandre Led og Udtryk, er at gaae til en
Yderlighed; at ville ved en splinterny Aabenbaring sikkre
sig en bogstavelig Gjengivelse af den sproglige Form, og
den udvortes Maade, paa hvilken Grundordet oprindelig er
meddelt, er at gaae til en modsat Yderlighed. I det
første Extrem gaaer man udenom Ordet for at gribe
Aanden; i det sidste Extrem gaaer man videre end Aanden
for at afstive Ordet. Sandheden er Extremernes
Midte, er Aandsvirkelighed, Virkelighed for Aanden.

\begin{center}VI.\end{center}

\emph{Aandsvirkelighed, Sandhedsmærket for
Troesbekjendelsen og den hellige Skrift, gjør
den historiske Troværdighed absolut afhængig
af den religiøse og giver Troeskritiken Maalestok
til Adskillelse af Aabenbaring og Mythologi}.

Saa stor end Afstanden er mellem det syttende Aarhundredes
skriftdogmatiske Aabenbaringslære og Nutidens
bibelkritiske Mythologi, saa er der dog een Grundsætning,
der er fælles for begge.

 »Enhver indre, af Tilegnelsen betinget Forandring,
som en udvortes Begivenhed undergaaer i den Fortællendes
Fremstilling, \emph{er en Forvanskning} af dens Virkelighed«;
dette er den øverste Præmis, den logiske Oversætning,
hvorefter den orthodoxe Apologetik og den
negative Kritik have formuleret hver sin Slutning. ---
»De hellige Skribenter have ikke ladet Begivenhederne
% --- p. 35 ---
\setcounter{footnote}{0}%
\opage{35}undergaae nogensomhelst Ændring; \emph{ergo}: er den hellige
Historie i alle Punkter og Dele en med sig selv
ordret overensstemmende, udvortes virkelig Historie«. Saaledes
slutter Orthodoxien. --- »De bibelske Forfattere have
i enhver Fremstilling af det Overnaturlige ladet Begivenheden
undergaae en, om end ubevidst, saa dog væsenlig
Forvandling; \emph{ergo}: er den hele saakaldte hellige Historie
en i alt Væsenligt forvansket, uvirkelig Historie«.
Saaledes slutter den negative Kritik. Mæglingstheologien
mægler og mægler, men uden Holdning, uden Princip.

Er da virkelig enhver ved troende Tilegnelse foregaaet
Forandring, som en overnaturlig Begivenhed kan
undergaae, en Forvanskning? Dette for Bibelens Troværdighed
afgjørende Grundspørgsmaal, maa besvares med
Ja eller Nei.

Besvares Spørgsmaalet med Ja, saa er Bibelordet
mythisk, og den opløsende Videnskritik faaer i Hovedsagen
Ret. »Muligheden til Mythe bliver til Sandsynlighed,
naar det Vidunderlige i Fortællingerne, de uopløselige,
kun ved den konstigste Harmonistik forklarlige
Modsigelser mellem de enkelte Evangelister, de mangehaande
kronologiske Vanskeligheder, der ikke lade sig
bringe i Overensstemmelse med samtidige Profanskribenter,
blive tagne med i Betragtning. Ei heller er det
tilstrækkeligt, som man hidtil brugte, at udligne alle disse
Modsigelser derved, at een Evangelist prisgives, en anden
begunstiges, saa at Matthæus bliver bragt som Offer for
Lukas, og denne, naar det fordres, atter for Johannes.
Dette er blot at maale med ulige Maal og Vægt. Det
viser sig tilmed ved en upartisk Betragtning af de forskjellige
historiske Fremstillinger, at alle Evangelister i
% --- p. 36 ---
\setcounter{footnote}{0}%
\opage{36}lige Grad maae underkjendes, at den enes Vidnesbyrd er
ligesaa meget, eller rettere sagt, ligesaa lidet værd som
den andens. Her føres en bellum omnium contra omnes,
og ingensteds formaae vi at vinde fast historisk Grund«\footnote[1]{K. Schwarz. Den nyeste Theologies Historie i kort Fremstilling (oversat efter Originalens tredie Oplag. Kjøbenhavn 1869). S. 94.}.

Besvares Spørgsmaalet med Nei, da træder Bibelordet
paany i Kraft som Aabenbaringsord; Troeslæren
har ved dette Nei gjort sig fri for Opløsningskritikens
og den mythiske Fortolknings Indsigelser og Paahæng.
Det ligger, som ovenfor viist, i Sagens Natur, at Under
og Tro maae i Opfattelse og Fremstilling være uadskillelige;
den Forvandling, Opfattelsen har medført, er
en Forvandling i Aanden. Forvandlingen er kun da en
Forvanskning, naar Aanden er falsk. Den mythiske Aand,
der sætter det Uvirkelige som virkeligt, er falsk Aand.
Men Aabenbaringens Aand er Sandhedens Aand, fordi
den i dybeste Forstand er Virkelighedens Aand.

At der, hvad det Udvortes og Enkelte angaaer, kan
gives Feiltagelser, som ikke i fjerneste Maade forvanske
Troesindholdet, vil kun en aandløs, bogstavbunden Supranaturalisme
kunne nægte. Der kan da ogsaa fra Evangelietroens
Synspunkt gjøres principiel Forskjel mellem Sagn,
sagnagtige Træk, og Myther. Hvorledes der kunde »komme
Sagn ud iblandt Brødrene«, og det endog Sagn, der vare
foranledigede ved uskyldig Misforstaaelse af Jesu Ord,
sees af John. 21, 23. Men slige Sagn indblande ikke
noget fra Aabenbaringstanken afvigende, mod dens Væsen
stridende Element, forvanske ikke Virkelighedsprincipet
og ere følgelig ikke Myther. I Matth. 27, 51 fortælles,
% --- p. 37 ---
\setcounter{footnote}{0}%
\opage{37}at, da Jesus »udgav Aanden paa Korset, da skjalv Jorden
og Klipperne revnede, og Gravene oplodes, og mange
af de hensovede Helliges Legemer opstode, og de gik ud
af Gravene efter hans Opstandelse, og kom ind i den
hellige Stad og aabenbaredes for Mange». Dette kan,
om man endelig vil, opfattes som et Sagn, men det kan
paa ingen Maade betragtes som en Mythe. Er Opstandelsen
af de Jerusalemske Døde en Mythe, saa er Aabenbaringsprincipet
fornægtet, saa er Lazari Opvækkelse og
Kristi Opstandelse ogsaa Myther, thi det Ene gaaer frem
af det samme Princip som det Andet\footnote[1]{At Forholdet dog ikke kan vendes om, saa at Kristi Opstandelse blev et Sagn, og kun Opstandelsen af de hensovede Helliges Legemer en Virkelighed, forstaaes af sig selv.}.

Det er en let Sag at faae de bibelske Fortællinger
til at gjælde for fuld Virkelighed, naar det Hele er gjort
med, at man uden videre tager Alt efter Bogstaven,
støtter sig til blind Autoritet og sætter Paastande istedenfor
Grunde. Hvad er lettere end at sige: »Det staaer i
Bibelen, det er historisk, og \emph{derfor} skulle vi troe det;
dersom vi ikke forud havde uigjendrivelige Beviser for at
den hellige Skrift er punktligere end noget juridisk Dokument,
saa havde vor Tro intet sikkert Grundlag; men det
staaer skrevet, \emph{det staaer skrevet, det staaer bogstavelig
skrevet}! --- Ved at aabne Modsigelsens Dør og
lyse den skriftjuridiske Paastaaelighed ned ad Trappen har
Videnskritiken, menneskelig talt, gjort Bibelordet en sand
Tjeneste.

Der er Virkelighed nok, »fuld Virkelighed« i de aabenbarede
Kjendsgjerninger; den hellige Historie er som et
rigt dækket Bord: træder ind, thi Alt er beredt. Men
% --- p. 38 ---
\setcounter{footnote}{0}%
\opage{38}det saaledes Beredte kommer kun dem tilgode, som værdigen
forstaae at nyde. At erklære de overnaturlige Næringsmidler
for Skueretter, fordi man, ved at indaande
Gift, har mistet Fordøielsen, er tragisk letfærdigt, frivolt.
At sætte sig til Bords og fylde sine dogmatiske Lommer
med udvortes Stof, som man agter at proppe i Andre,
men ikke selv kan fordøie, er en plump Misbrug af det
Hellige. Aabenbaringens Virkelighed er kun Virkelighed
for dem, som i Hjertets Enfoldighed have Tilegnelsens
Trang og i Troens Forstand Tilegnelsens Gave. Af anden
Virkelighed kan man efter Omstændighederne snart faae
for Lidt, snart igjen for Meget, desværre alt for Meget;
men af Aabenbaringens Virkelighed kan man i Troen altid
være vis paa hverken at faae for Lidt eller for Meget;
thi Tilegnelsen er Maalet, og Tilegnelsen betinges af
Trangen. Aabenbaringsvirkelighed er i det Udvortes som
i det Indvortes Aandsvirkelighed, en Virkelighed, som paa
Tilegnelsens og Selvforstaaelsens Vilkaar fra første Begyndelse
af har været, er og i al Evighed vil være Virkelighed
for Aanden.
\signature{R. \emph{Nielsen}.}

\end{document}
